% Overleaf: ustaw Compiler na pdfLaTeX. Plik nie wymaga lokalnych fontów.
\documentclass[11pt,a4paper]{article}

\usepackage{iftex}
\ifPDFTeX
  \usepackage[utf8]{inputenc}
  \usepackage[T1]{fontenc}
  \usepackage[polish]{babel}
  \usepackage{lmodern}
\else
  \usepackage{fontspec}
  \setmainfont{Times New Roman}
  \setsansfont{Arial}
  \setmonofont{Courier New}
  \usepackage[polish]{babel}
\fi

\usepackage[a4paper,margin=2.35cm]{geometry}
\usepackage{amsmath,amssymb,mathtools}
\usepackage{booktabs,longtable,array}
\usepackage{enumitem}
\usepackage{hyperref}
\usepackage{xcolor}
\usepackage{listings}
\usepackage{fancyhdr}
\usepackage{titlesec}
\usepackage{tocloft}

\hypersetup{
  colorlinks=true,
  linkcolor=blue!50!black,
  urlcolor=blue!60!black,
  citecolor=blue!50!black,
  pdftitle={Opracowanie pytań na egzamin dyplomowy MSI},
  pdfauthor={opracowanie wygenerowane z materiałów lokalnych}
}

\renewcommand{\contentsname}{Spis treści}
\renewcommand{\figurename}{Rysunek}
\renewcommand{\tablename}{Tabela}
\renewcommand{\refname}{Źródła}
\setlength{\parindent}{0pt}
\setlength{\parskip}{0.55em}
\setlist[itemize]{leftmargin=1.4em,itemsep=0.18em,topsep=0.25em}
\setlist[enumerate]{leftmargin=1.6em,itemsep=0.18em,topsep=0.25em}
\setlist[description]{leftmargin=1.8em,itemsep=0.18em,topsep=0.25em}
\sloppy

\pagestyle{fancy}
\fancyhf{}
\lhead{Egzamin dyplomowy MSI}
\rhead{\thepage}
\renewcommand{\headrulewidth}{0.4pt}

\titleformat{\section}{\Large\bfseries\sffamily}{\thesection}{0.8em}{}
\titleformat{\subsection}{\large\bfseries\sffamily}{\thesubsection}{0.8em}{}
\titleformat{\subsubsection}{\normalsize\bfseries\sffamily}{\thesubsubsection}{0.8em}{}

\newcommand{\R}{\mathbb{R}}
\newcommand{\N}{\mathbb{N}}
\newcommand{\BigO}{\mathcal{O}}
\newcommand{\NP}{\mathrm{NP}}
\newcommand{\coNP}{\mathrm{coNP}}
\newcommand{\SAT}{\mathrm{SAT}}
\newcommand{\argmax}{\operatorname*{arg\,max}}
\newcommand{\argmin}{\operatorname*{arg\,min}}
\newcommand{\eps}{\varepsilon}
\newcommand{\todo}[1]{\textcolor{red!70!black}{#1}}

\lstset{
  basicstyle=\ttfamily\small,
  breaklines=true,
  columns=fullflexible,
  keepspaces=true,
  frame=single,
  framerule=0.2pt,
  xleftmargin=0.5em,
  xrightmargin=0.5em
}

\begin{document}

\begin{titlepage}
  \centering
  {\Huge\bfseries Bardzo dokładne opracowanie pytań na egzamin dyplomowy\par}
  \vspace{0.6cm}
  {\Large Informatyka i Systemy Informacyjne, specjalność Metody Sztucznej Inteligencji\par}
  \vspace{1.2cm}
  {\large Opracowanie przygotowane na podstawie listy pytań oraz materiałów w folderze projektu.\par}
  \vfill
  {\large 24 czerwca 2026\par}
\end{titlepage}

\tableofcontents
\newpage

\section*{Jak korzystać z opracowania}
\addcontentsline{toc}{section}{Jak korzystać z opracowania}

Dokument jest ułożony według pytań z listy, a nie według plików z wykładami. Przy każdym zagadnieniu najpierw podane są definicje i intuicja, potem formalizacja, typowe algorytmy albo twierdzenia, a na końcu elementy, które warto umieć powiedzieć ustnie na obronie.

Materiały lokalne wykorzystane bezpośrednio:
\begin{itemize}
  \item \texttt{Zalacznik-3-Lista-przykladowych-pytan-ISI-MSI-2020-12-16 copy.pdf} - lista pytań.
  \item \texttt{Slajdy/RW\_01--RW\_14.pdf} - reprezentacja wiedzy, rezolucja, wnioskowanie o działaniach.
  \item \texttt{merged.pdf} - uczenie ze wzmocnieniem.
  \item \texttt{wykłady\_m\_scalone.pdf} - sieci neuronowe.
\end{itemize}

\section*{Lista pytań objętych dokumentem}
\addcontentsline{toc}{section}{Lista pytań objętych dokumentem}

Opracowanie obejmuje wszystkie zagadnienia z listy: hierarchię Chomsky'ego, złożoność obliczeniową, klasy problemów i P--NP, testowanie metod inteligencji obliczeniowej, eksperymenty obliczeniowe, TSP, interpolację, metody liniowe, zera funkcji, całkowanie numeryczne, programowanie matematyczne, algorytmy aproksymacyjne, programowanie dynamiczne, zarządzanie projektami IT, rezolucję, wnioskowanie o działaniach, algorytmy ewolucyjne, Monte Carlo w grach, agentów, inteligencję rojową, Prolog, reguły minimalne, redukcję wymiaru, selekcję cech, braki danych, estymatory błędów, uczenie ze wzmocnieniem, MDP, TD-learning oraz sieci neuronowe.

\section{Podstawy teoretyczne, algorytmiczne i numeryczne}

\subsection{Hierarchia Chomsky'ego oraz odpowiadające jej gramatyki i automaty}

Hierarchia Chomsky'ego porządkuje klasy języków formalnych według tego, jak skomplikowanych reguł gramatycznych potrzeba do ich opisu oraz jak silny model automatu jest potrzebny do ich rozpoznawania. Im wyższy poziom w hierarchii, tym większa ekspresyjność, ale też trudniejsze własności decyzyjne.

\subsubsection{Podstawowe pojęcia}

Alfabet $\Sigma$ jest skończonym zbiorem symboli. Słowo to skończony ciąg symboli z $\Sigma$, a język to dowolny podzbiór $\Sigma^*$. Gramatyka formalna jest zwykle zapisywana jako
\[
G=(N,\Sigma,P,S),
\]
gdzie $N$ jest zbiorem symboli nieterminalnych, $\Sigma$ zbiorem terminali, $P$ zbiorem produkcji, a $S\in N$ symbolem startowym. Produkcja mówi, jak wolno zastąpić fragment sentencji innym fragmentem. Język generowany przez gramatykę to zbiór słów terminalnych wyprowadzalnych z $S$.

Automat jest modelem rozpoznawania języka. Dla słowa wejściowego automat odpowiada, czy słowo należy do języka. Różne klasy automatów mają różną pamięć i różne możliwości sterowania obliczeniem.

\subsubsection{Cztery klasy w hierarchii}

\begin{longtable}{p{0.13\textwidth}p{0.23\textwidth}p{0.26\textwidth}p{0.25\textwidth}}
\toprule
Typ & Gramatyka & Ograniczenie produkcji & Odpowiadający automat \\
\midrule
0 & bez ograniczeń & $\alpha\to\beta$, gdzie $\alpha$ zawiera co najmniej jeden nieterminal & maszyna Turinga \\
1 & kontekstowa & produkcje niemalejące długości, często $\alpha A\beta\to\alpha\gamma\beta$, $\gamma\ne\epsilon$ & liniowo ograniczona maszyna Turinga \\
2 & bezkontekstowa & $A\to\gamma$, po lewej dokładnie jeden nieterminal & automat ze stosem \\
3 & regularna & np. $A\to aB$, $A\to a$, ewentualnie $A\to\epsilon$ w ograniczonej postaci & automat skończony \\
\bottomrule
\end{longtable}

Zachodzą ścisłe inkluzje:
\[
\mathcal{L}_{3}\subsetneq \mathcal{L}_{2}\subsetneq \mathcal{L}_{1}\subsetneq \mathcal{L}_{0}.
\]
Każdy język regularny jest bezkontekstowy, każdy bezkontekstowy jest kontekstowy, a każdy kontekstowy jest rekurencyjnie przeliczalny. Odwrotne zależności nie zachodzą.

\subsubsection{Języki regularne}

Języki regularne opisują najprostsze wzorce, bez pamięci o nieograniczonej strukturze zagnieżdżonej. Są równoważne:
\begin{itemize}
  \item gramatykom regularnym,
  \item wyrażeniom regularnym,
  \item deterministycznym automatom skończonym,
  \item niedeterministycznym automatom skończonym.
\end{itemize}

Automat skończony ma skończoną liczbę stanów i nie ma stosu ani taśmy roboczej. Dlatego potrafi zapamiętać tylko skończoną informację o prefiksie wejścia. Typowy język regularny:
\[
L=\{w\in\{0,1\}^*: w\text{ kończy się na }01\}.
\]
Typowy język nieregularny:
\[
L=\{a^n b^n:n\ge 0\}.
\]
Nie jest regularny, bo wymaga porównania dwóch nieograniczonych liczności.

\subsubsection{Języki bezkontekstowe}

Języki bezkontekstowe pozwalają opisywać strukturę rekurencyjną, np. nawiasowanie, drzewa składniowe i większość składni języków programowania. Produkcja ma postać $A\to\gamma$, więc nieterminal może być rozwijany niezależnie od kontekstu.

Automat ze stosem ma dodatkową pamięć LIFO. Dzięki temu rozpoznaje np.
\[
\{a^n b^n:n\ge 0\}
\]
przez odkładanie symboli dla liter $a$ i zdejmowanie ich przy literach $b$. Nie rozpoznaje jednak w ogólności języków wymagających dwóch niezależnych liczników, np.
\[
\{a^n b^n c^n:n\ge 0\}.
\]

\subsubsection{Języki kontekstowe}

Gramatyki kontekstowe pozwalają na produkcje zależne od otoczenia symbolu. Warunek niemalejącej długości powoduje, że wyprowadzenie nie może dowolnie skracać sentencji. Odpowiada im liniowo ograniczona maszyna Turinga, czyli maszyna Turinga, która może korzystać tylko z pamięci liniowo ograniczonej rozmiarem wejścia. Przykładem języka kontekstowego niebędącego bezkontekstowym jest
\[
\{a^n b^n c^n:n\ge 1\}.
\]

\subsubsection{Języki typu 0}

Gramatyki typu 0 nie mają istotnych ograniczeń na produkcje poza tym, że lewa strona nie może być pusta. Odpowiadają maszynom Turinga i językom rekurencyjnie przeliczalnym: jeśli słowo należy do języka, maszyna kiedyś je zaakceptuje; jeśli nie należy, może odrzucić albo działać w nieskończoność. Ta klasa obejmuje więc problemy częściowo rozstrzygalne, ale niekoniecznie decyzyjne.

\subsubsection{Co powiedzieć na obronie}

Najważniejsza odpowiedź ustna: hierarchia Chomsky'ego klasyfikuje języki według ograniczeń reguł produkcji. Typ 3 to języki regularne i automaty skończone, typ 2 to języki bezkontekstowe i automaty ze stosem, typ 1 to języki kontekstowe i liniowo ograniczone maszyny Turinga, typ 0 to gramatyki bez ograniczeń i maszyny Turinga. Każda kolejna klasa jest ściśle szersza.

\subsection{Złożoność obliczeniowa algorytmu}

Złożoność obliczeniowa opisuje, jak zasoby potrzebne algorytmowi rosną wraz z rozmiarem danych wejściowych. Najczęściej analizuje się czas i pamięć, ale w praktyce mogą liczyć się też liczba zapytań do bazy, komunikacja sieciowa, zużycie energii albo liczba operacji arytmetycznych wysokiej precyzji.

\subsubsection{Rozmiar wejścia i model obliczeń}

Złożoność zawsze definiuje się względem parametru $n$, czyli rozmiaru wejścia. Dla tablicy może to być liczba elementów, dla grafu para $(|V|,|E|)$, dla liczby całkowitej liczba bitów jej zapisu, a nie sama wartość liczby. To rozróżnienie jest ważne: algorytm wielomianowy w wartości liczby może być wykładniczy w długości jej zapisu.

Model obliczeń określa, jakie operacje uznajemy za elementarne. W klasycznej analizie RAM zakłada się, że proste operacje arytmetyczne, porównania i dostęp do pamięci zajmują czas stały. W analizie bitowej koszt operacji zależy od długości operandów.

\subsubsection{Notacja asymptotyczna}

Dla funkcji kosztu $T(n)$ używa się:
\begin{itemize}
  \item $T(n)=\BigO(g(n))$ - ograniczenie z góry z dokładnością do stałej,
  \item $T(n)=\Omega(g(n))$ - ograniczenie z dołu,
  \item $T(n)=\Theta(g(n))$ - dokładny rząd wzrostu,
  \item $T(n)=o(g(n))$ - wzrost istotnie wolniejszy,
  \item $T(n)=\omega(g(n))$ - wzrost istotnie szybszy.
\end{itemize}

Formalnie $T(n)=\BigO(g(n))$, gdy istnieją stałe $c>0$ i $n_0$, takie że dla każdego $n\ge n_0$ zachodzi $0\le T(n)\le c g(n)$.

\subsubsection{Przypadek pesymistyczny, średni i amortyzowany}

\begin{description}
  \item[Pesymistyczna złożoność] mierzy najgorszy możliwy koszt dla wejść rozmiaru $n$. Daje gwarancję niezależną od danych.
  \item[Średnia złożoność] wymaga rozkładu prawdopodobieństwa na wejściach. Jest użyteczna, ale zależy od założeń statystycznych.
  \item[Złożoność oczekiwana] dotyczy zwykle algorytmów losowych i oczekiwania po losowości algorytmu.
  \item[Złożoność amortyzowana] rozlicza koszt sekwencji operacji. Pojedyncza operacja może być droga, ale średni koszt w długiej sekwencji jest mały, np. dynamiczna tablica z podwajaniem rozmiaru.
\end{description}

\subsubsection{Typowe rzędy wzrostu}

Od najlepszych do najgorszych praktycznie:
\[
\BigO(1),\quad \BigO(\log n),\quad \BigO(n),\quad \BigO(n\log n),\quad
\BigO(n^2),\quad \BigO(n^3),\quad \BigO(2^n),\quad \BigO(n!),\quad \BigO(n^n).
\]
Różnice asymptotyczne szybko dominują nad stałymi. Algorytm $\BigO(n\log n)$ dla dużych danych będzie zwykle jakościowo lepszy od $\BigO(n^2)$, nawet jeśli ma większą stałą.

\subsubsection{Złożoność pamięciowa}

Złożoność pamięciowa $S(n)$ mierzy dodatkową pamięć potrzebną algorytmowi. Często odróżnia się pamięć wejścia od pamięci roboczej. Algorytm sortowania przez scalanie ma czas $\Theta(n\log n)$ i typowo pamięć $\Theta(n)$, a heapsort czas $\Theta(n\log n)$ i pamięć dodatkową $\Theta(1)$.

\subsubsection{Co powiedzieć na obronie}

Złożoność algorytmu to funkcja opisująca zużycie zasobów w zależności od rozmiaru wejścia. Podajemy ją asymptotycznie, najczęściej w notacji $\BigO$, bo interesuje nas tempo wzrostu dla dużych danych. Trzeba wskazać model obliczeń, rozmiar wejścia i to, czy mówimy o czasie, pamięci, przypadku pesymistycznym, średnim czy amortyzowanym.

\subsection{Klasy problemów według złożoności. Problem P--NP}

Klasy złożoności grupują problemy decyzyjne według zasobów potrzebnych do ich rozwiązania. Problem decyzyjny ma odpowiedź TAK albo NIE. Wiele problemów optymalizacyjnych można badać przez wersje decyzyjne, np. zamiast "znajdź najkrótszą trasę TSP" pytamy "czy istnieje trasa o długości co najwyżej $B$?".

\subsubsection{Klasa P}

Klasa $\mathrm{P}$ zawiera problemy decyzyjne rozwiązywalne deterministycznie w czasie wielomianowym względem długości wejścia:
\[
\mathrm{P}=\bigcup_{k\ge 1}\mathrm{DTIME}(n^k).
\]
Przykłady: osiągalność w grafie, najkrótsze ścieżki z nieujemnymi wagami w wersji decyzyjnej, sortowanie jako problem funkcyjny, sprawdzanie spójności grafu, maksymalny przepływ w wersji decyzyjnej.

Wielomianowość traktuje się jako formalny odpowiednik obliczalności efektywnej, choć praktycznie algorytm $\BigO(n^{100})$ nie jest użyteczny, a algorytm wykładniczy może działać dla małych $n$.

\subsubsection{Klasa NP}

Klasa $\NP$ zawiera problemy decyzyjne, dla których odpowiedź TAK ma certyfikat weryfikowalny deterministycznie w czasie wielomianowym. Równoważnie są to problemy rozwiązywalne przez niedeterministyczną maszynę Turinga w czasie wielomianowym.

Przykład: dla problemu Hamiltona certyfikatem jest kolejność wierzchołków. Weryfikacja, czy jest to cykl Hamiltona, jest wielomianowa. Dla TSP decyzyjnego certyfikatem jest permutacja miast; weryfikujemy, czy tworzy cykl i czy suma wag jest $\le B$.

Zawsze $\mathrm{P}\subseteq\NP$, bo jeśli umiemy problem szybko rozwiązać, to umiemy też szybko zweryfikować certyfikat albo go zignorować.

\subsubsection{Redukcje wielomianowe}

Problem $A$ redukuje się wielomianowo do problemu $B$, co zapisujemy $A\le_p B$, jeśli istnieje funkcja obliczalna w czasie wielomianowym taka, że instancja $x$ problemu $A$ ma odpowiedź TAK wtedy i tylko wtedy, gdy $f(x)$ ma odpowiedź TAK w problemie $B$.

Redukcja oznacza: gdybyśmy umieli szybko rozwiązywać $B$, umielibyśmy szybko rozwiązywać $A$. Dlatego trudność przenosi się z $A$ na $B$.

\subsubsection{Problemy NP-trudne i NP-zupełne}

Problem jest NP-trudny, jeśli każdy problem z $\NP$ redukuje się do niego w czasie wielomianowym. Problem jest NP-zupełny, jeśli jest NP-trudny i jednocześnie należy do $\NP$.

Klasyczne problemy NP-zupełne:
\begin{itemize}
  \item $\SAT$ i 3-SAT,
  \item CLIQUE w wersji decyzyjnej,
  \item VERTEX COVER,
  \item HAMILTONIAN CYCLE,
  \item TSP decyzyjny,
  \item SUBSET SUM.
\end{itemize}

Twierdzenie Cooka-Levina mówi, że $\SAT$ jest NP-zupełny. W praktyce wiele dowodów NP-zupełności polega na redukcji z problemu już znanego jako NP-zupełny.

\subsubsection{Problem P--NP}

Pytanie $\mathrm{P}=\NP?$ brzmi: czy każdy problem, którego rozwiązanie można szybko zweryfikować, można też szybko rozwiązać? Wiadomo, że $\mathrm{P}\subseteq\NP$, ale nie wiadomo, czy inkluzja jest ścisła. Powszechne przekonanie brzmi $\mathrm{P}\ne\NP$, ponieważ przez dekady nie znaleziono algorytmów wielomianowych dla problemów NP-zupełnych mimo ogromnego wysiłku.

Jeżeli $\mathrm{P}=\NP$, wszystkie problemy NP-zupełne miałyby algorytmy wielomianowe. Miałoby to konsekwencje dla optymalizacji, automatycznego dowodzenia, kryptografii i wielu problemów planowania. Jeżeli $\mathrm{P}\ne\NP$, istnieją problemy łatwo weryfikowalne, ale niemożliwe do rozwiązania w czasie wielomianowym przez algorytm deterministyczny.

\subsubsection{Inne klasy warte wspomnienia}

\begin{description}
  \item[co-NP] problemy, których odpowiedzi NIE mają krótkie certyfikaty. Nie wiadomo, czy $\NP=\coNP$.
  \item[PSPACE] problemy rozwiązywalne w pamięci wielomianowej. Obejmuje $\NP$.
  \item[EXPTIME] problemy rozwiązywalne w czasie wykładniczym.
  \item[BPP] problemy rozwiązywalne przez algorytmy losowe w czasie wielomianowym z ograniczonym prawdopodobieństwem błędu.
  \item[RP, co-RP, ZPP] klasy losowe z jednostronnym błędem albo zerowym błędem oczekiwanym.
\end{description}

\subsubsection{Co powiedzieć na obronie}

P to problemy decyzyjne rozwiązywalne w czasie wielomianowym. NP to problemy, dla których rozwiązanie można zweryfikować w czasie wielomianowym. Problem P--NP pyta, czy szybka weryfikacja oznacza szybkie rozwiązywanie. Problemy NP-zupełne są najtrudniejsze w NP: jeśli jeden z nich ma algorytm wielomianowy, to $\mathrm{P}=\NP$.

\subsection{Testowanie i poszukiwanie błędów w implementacji metod inteligencji obliczeniowej wykorzystujących losowość}

Metody inteligencji obliczeniowej często są niedeterministyczne: algorytmy ewolucyjne, symulowane wyżarzanie, Monte Carlo Tree Search, losowe inicjalizacje sieci neuronowych, dropout, losowe mini-batche, bagging albo metody próbkowania. Testowanie takich implementacji wymaga odróżnienia błędów programistycznych od naturalnej zmienności wyników.

\subsubsection{Reprodukowalność jako pierwszy warunek debugowania}

Podstawą jest kontrola generatorów liczb pseudolosowych. Należy:
\begin{itemize}
  \item jawnie ustawiać ziarno generatora,
  \item zapisywać użyte ziarna razem z wynikami eksperymentu,
  \item kontrolować wszystkie źródła losowości, np. bibliotekę numeryczną, framework GPU, operacje wielowątkowe,
  \item umożliwić powtórzenie pojedynczego przebiegu bit po bicie, o ile środowisko na to pozwala.
\end{itemize}

Dobrym wzorcem jest rozdzielenie generatorów: osobny strumień dla inicjalizacji populacji, mutacji, selekcji, podziału danych, szumu eksploracyjnego. Ułatwia to znalezienie źródła rozbieżności.

\subsubsection{Testy jednostkowe deterministycznych komponentów}

Nawet metoda losowa składa się z wielu deterministycznych części. Testować należy osobno:
\begin{itemize}
  \item funkcję celu i ograniczenia,
  \item kod dekodowania reprezentacji osobnika,
  \item obliczanie dopasowania,
  \item operatory krzyżowania i mutacji dla ustalonych wartości losowych,
  \item aktualizację wag, gradientów, wartości $Q$ albo statystyk w drzewie MCTS,
  \item obsługę przypadków brzegowych: puste dane, jeden osobnik, zerowa wariancja, remisy, wartości NaN.
\end{itemize}

Warto wstrzykiwać generator losowy jako zależność. Wtedy test może użyć generatora deterministycznego zwracającego z góry znaną sekwencję.

\subsubsection{Testy własności i niezmienników}

W metodach inteligencji obliczeniowej często łatwiej sprawdzić własności niż konkretne wyniki:
\begin{itemize}
  \item prawdopodobieństwa selekcji sumują się do $1$,
  \item mutacja zachowuje dopuszczalny zakres genów,
  \item krzyżowanie permutacyjne zwraca permutację bez duplikatów,
  \item najlepsze znalezione rozwiązanie w algorytmie elitarnym nie pogarsza się między generacjami,
  \item operator projekcji zwraca punkt spełniający ograniczenia,
  \item funkcja straty maleje dla prostego problemu uczącego, jeżeli krok jest odpowiednio mały,
  \item estymator Monte Carlo jest nieobciążony na prostym rozkładzie testowym.
\end{itemize}

Takie testy można realizować jako property-based testing: generujemy wiele losowych wejść i sprawdzamy ogólne warunki.

\subsubsection{Testy statystyczne losowości}

Jeżeli komponent ma realizować określony rozkład, pojedynczy wynik niczego nie dowodzi. Sprawdza się zbiorcze statystyki:
\begin{itemize}
  \item średnią i wariancję próbek,
  \item histogram kategorii,
  \item test $\chi^2$ dla rozkładu dyskretnego,
  \item test Kołmogorowa-Smirnowa dla rozkładu ciągłego,
  \item przedziały ufności dla oczekiwanej jakości.
\end{itemize}

Przykład: jeśli mutacja wybiera każdy gen z prawdopodobieństwem $p$, to po wielu próbach liczba mutacji powinna mieścić się w sensownym przedziale wokół $np$. Test nie powinien być zbyt wąski, bo będzie fałszywie wykrywać błędy.

\subsubsection{Metamorphic testing}

Metamorphic testing polega na sprawdzaniu relacji między wynikami dla przekształconych danych wejściowych. Przykłady:
\begin{itemize}
  \item przemnożenie funkcji celu przez dodatnią stałą nie powinno zmienić optimum,
  \item permutacja kolejności przykładów uczących nie powinna zmienić wyniku batchowego,
  \item dodanie stałej do wszystkich długości krawędzi w TSP może zmienić wartości tras, ale nie musi zmienić relacji, jeśli dodatek jest równy dla każdej trasy,
  \item powielenie identycznych próbek nie powinno zmieniać optimum empirycznego dla średniej straty.
\end{itemize}

\subsubsection{Debugowanie przebiegów optymalizacji}

Należy logować nie tylko wynik końcowy, ale też przebieg:
\begin{itemize}
  \item najlepszą, średnią i najgorszą wartość funkcji celu w populacji,
  \item różnorodność populacji,
  \item rozkład długości kroków albo mutacji,
  \item temperaturę w symulowanym wyżarzaniu,
  \item stopę akceptacji ruchów,
  \item normy gradientów i wartości straty,
  \item liczbę wizyt i wartości w węzłach MCTS.
\end{itemize}

Typowe symptomy błędów: natychmiastowa utrata różnorodności, wartości NaN, brak jakiejkolwiek zmiany mimo mutacji, nierealnie szybka zbieżność, wyniki lepsze niż znane dolne ograniczenie, silna zależność od kolejności danych, brak powtarzalności mimo ustalonego ziarna.

\subsubsection{Porównanie z bazami odniesienia}

Implementację sprawdza się na:
\begin{itemize}
  \item instancjach małych, dla których da się policzyć optimum brutalnie,
  \item problemach analitycznych z oczywistym optimum, np. funkcja sferyczna $f(x)=\sum_i x_i^2$,
  \item benchmarkach literaturowych,
  \item baseline'ach prostych, np. losowe przeszukiwanie, algorytm zachłanny, metoda lokalna.
\end{itemize}

Jeżeli złożona metoda regularnie przegrywa z losowym baseline'em na prostym zadaniu, zwykle oznacza to błąd implementacyjny, złe skalowanie parametrów albo niedopasowanie reprezentacji.

\subsubsection{Co powiedzieć na obronie}

Testując metody losowe, trzeba zapewnić reprodukowalność przez kontrolę ziaren, testować deterministyczne komponenty, sprawdzać niezmienniki i własności statystyczne, uruchamiać wiele niezależnych przebiegów oraz porównywać z prostymi benchmarkami. Wynik pojedynczego uruchomienia nie jest dowodem poprawności ani jakości.

\subsection{Metody prowadzenia eksperymentów obliczeniowych}

Eksperyment obliczeniowy ma odpowiedzieć na pytanie badawcze, a nie tylko wygenerować tabelę wyników. Wymaga hipotezy, kontrolowanego protokołu, rzetelnego doboru danych, powtarzalności i analizy statystycznej.

\subsubsection{Pytanie badawcze i hipotezy}

Najpierw definiuje się, co porównujemy i dlaczego. Przykłady:
\begin{itemize}
  \item Czy metoda A daje mniejszy błąd klasyfikacji niż metoda B?
  \item Czy operator mutacji X poprawia jakość algorytmu ewolucyjnego na instancjach TSP?
  \item Jak rozmiar populacji wpływa na czas dojścia do rozwiązania?
  \item Czy metoda jest odporna na braki danych?
\end{itemize}

Hipoteza zerowa może mówić, że nie ma różnicy między metodami. Hipoteza alternatywna mówi, że różnica istnieje, np. metoda A ma mniejszą średnią stratę.

\subsubsection{Plan eksperymentu}

Dobry plan określa:
\begin{itemize}
  \item dane lub instancje problemu,
  \item metody i ich implementacje,
  \item parametry i sposób ich strojenia,
  \item metryki jakości,
  \item budżet obliczeniowy,
  \item liczbę powtórzeń,
  \item sposób agregacji wyników,
  \item testy statystyczne,
  \item kryteria zakończenia.
\end{itemize}

Ważna zasada: zbiór testowy nie może służyć do strojenia parametrów. W uczeniu maszynowym typowy podział to train/validation/test albo walidacja krzyżowa z zewnętrzną pętlą testową.

\subsubsection{Kontrola zmiennych}

Aby porównanie było uczciwe, zmieniamy jedną rzecz naraz albo stosujemy formalny plan czynnikowy. Wszystkie metody powinny mieć porównywalny budżet: czas, liczba ewaluacji funkcji celu, liczba epok, liczba symulacji. W algorytmach optymalizacyjnych często bardziej uczciwe jest porównanie po liczbie wywołań funkcji celu niż po liczbie iteracji, bo jedna iteracja różnych metod może mieć różny koszt.

\subsubsection{Powtarzalność i replikowalność}

Powtarzalność oznacza możliwość uzyskania tych samych wyników w tym samym środowisku. Replikowalność oznacza potwierdzenie w innym środowisku albo przez inną implementację. Należy zapisywać:
\begin{itemize}
  \item kod i wersję commita,
  \item wersje bibliotek,
  \item konfigurację sprzętu,
  \item parametry eksperymentu,
  \item ziarna losowości,
  \item surowe wyniki,
  \item skrypty do agregacji.
\end{itemize}

\subsubsection{Metryki i analiza statystyczna}

Metryka musi odpowiadać celowi. Dla klasyfikacji używa się accuracy, precision, recall, F1, AUC, log-loss; dla regresji MAE, MSE, RMSE, $R^2$; dla optymalizacji wartości funkcji celu, odchylenia od optimum, czasu, liczby ewaluacji, stabilności.

Dla metod losowych podaje się co najmniej średnią i odchylenie standardowe albo medianę i kwartyle. Przy rozkładach skośnych mediana bywa bardziej informacyjna niż średnia. Do porównań można użyć testu t dla sparowanych wyników, testu Wilcoxona, testu Friedmana z post-hoc dla wielu metod i wielu zbiorów danych. Sama różnica średnich bez niepewności jest słabym argumentem.

\subsubsection{Ablation study i analiza czułości}

Ablation study polega na usuwaniu albo podmienianiu składników metody, aby sprawdzić, które elementy naprawdę wpływają na wynik. Analiza czułości bada, jak wynik zmienia się przy zmianie parametrów. To ważne, bo metoda może działać dobrze tylko dla wąskiego zakresu parametrów, co zmniejsza jej praktyczną użyteczność.

\subsubsection{Błędy w eksperymentach}

Typowe błędy:
\begin{itemize}
  \item przeciek danych między treningiem i testem,
  \item strojenie parametrów na zbiorze testowym,
  \item porównanie do słabego baseline'u,
  \item zbyt mała liczba powtórzeń,
  \item ignorowanie wariancji,
  \item wybiórcze raportowanie najlepszych wyników,
  \item inny budżet obliczeniowy dla różnych metod,
  \item brak kontroli losowości,
  \item wyciąganie wniosków z jednego datasetu albo jednej instancji.
\end{itemize}

\subsubsection{Co powiedzieć na obronie}

Eksperyment obliczeniowy powinien mieć jasną hipotezę, ustalony protokół, kontrolowane parametry, uczciwe baseline'y, powtórzenia, analizę statystyczną i możliwość odtworzenia. W metodach losowych trzeba raportować rozkład wyników, a nie tylko najlepszy przebieg.

\subsection{Problem komiwojażera: definicja i algorytmy rozwiązywania}

Problem komiwojażera, czyli TSP, jest jednym z najważniejszych problemów optymalizacji kombinatorycznej. Dane jest $n$ miast i koszty przejazdu $c_{ij}$ między parami miast. Należy znaleźć cykl Hamiltona o minimalnym koszcie, czyli trasę odwiedzającą każde miasto dokładnie raz i wracającą do punktu startowego.

\subsubsection{Warianty problemu}

\begin{description}
  \item[Symetryczny TSP] $c_{ij}=c_{ji}$. Trasa nie zależy od kierunku przejścia krawędzi.
  \item[Asymetryczny TSP] koszty mogą zależeć od kierunku, np. czas jazdy w mieście.
  \item[Metryczny TSP] koszty spełniają nierówność trójkąta $c_{ij}\le c_{ik}+c_{kj}$.
  \item[Euklidesowy TSP] miasta są punktami w przestrzeni, a kosztem jest odległość euklidesowa.
  \item[Decyzyjny TSP] pytanie brzmi, czy istnieje cykl o koszcie co najwyżej $B$.
\end{description}

Decyzyjna wersja TSP jest NP-zupełna, a wersja optymalizacyjna jest NP-trudna. Oznacza to, że nie oczekujemy algorytmu dokładnego wielomianowego dla ogólnego TSP, o ile $\mathrm{P}\ne\NP$.

\subsubsection{Sformułowanie matematyczne}

Dla symetrycznego TSP można wprowadzić zmienne binarne $x_{ij}$, gdzie $x_{ij}=1$, jeśli krawędź $(i,j)$ należy do trasy. Minimalizujemy
\[
\min \sum_{i<j} c_{ij}x_{ij}.
\]
Warunki stopnia:
\[
\sum_{j\ne i} x_{ij}=2\quad\text{dla każdego miasta }i.
\]
Same warunki stopnia dopuszczają jednak kilka rozłącznych cykli. Dlatego trzeba dodać ograniczenia eliminujące podcykle:
\[
\sum_{i,j\in S,\,i<j}x_{ij}\le |S|-1
\quad\text{dla każdego } \emptyset\ne S\subsetneq V.
\]
Tych ograniczeń jest wykładniczo wiele, ale w metodach branch-and-cut dodaje się tylko te naruszone.

\subsubsection{Algorytmy dokładne}

\begin{description}
  \item[Brute force] sprawdza wszystkie permutacje miast. Dla TSP symetrycznego liczba różnych tras to $(n-1)!/2$. Koszt jest silnie wykładniczy.
  \item[Programowanie dynamiczne Held-Karpa] definiuje $D[S,j]$ jako koszt najkrótszej ścieżki zaczynającej w mieście $1$, odwiedzającej zbiór $S$ i kończącej w $j$. Rekurencja:
  \[
  D[S,j]=\min_{i\in S,\,i\ne j}\left(D[S\setminus\{j\},i]+c_{ij}\right).
  \]
  Złożoność czasu $\BigO(n^2 2^n)$ i pamięci $\BigO(n2^n)$.
  \item[Branch and bound] przeszukuje drzewo częściowych tras i odcina gałęzie, których dolne ograniczenie jest gorsze od najlepszego znanego rozwiązania.
  \item[Branch and cut] łączy programowanie całkowitoliczbowe z generowaniem ograniczeń eliminujących podcykle i nierówności wzmacniających relaksację.
\end{description}

\subsubsection{Heurystyki konstrukcyjne}

Heurystyki szybko budują rozwiązanie, ale bez gwarancji optymalności:
\begin{itemize}
  \item najbliższy sąsiad: z aktualnego miasta idziemy do najbliższego nieodwiedzonego,
  \item nearest insertion: wstawiamy kolejne miasto w miejsce najmniej zwiększające koszt,
  \item cheapest insertion: wybieramy globalnie najtańsze wstawienie,
  \item farthest insertion: najpierw wstawiamy miasta dalekie, żeby wcześniej ustalić strukturę trasy,
  \item algorytmy oparte na minimalnym drzewie rozpinającym.
\end{itemize}

\subsubsection{Ulepszanie lokalne}

Metody lokalne startują z dowolnej trasy i poprawiają ją przez ruchy lokalne:
\begin{description}
  \item[2-opt] usuwa dwie krawędzie i odwraca fragment trasy, jeśli to skraca cykl.
  \item[3-opt] usuwa trzy krawędzie i rozważa kilka sposobów ponownego połączenia fragmentów.
  \item[Lin-Kernighan] adaptacyjnie dobiera złożoność ruchu $k$-opt i jest bardzo skuteczny praktycznie.
\end{description}

2-opt jest prosty i często daje dużą poprawę. Warunek poprawy dla wymiany krawędzi $(a,b)$ i $(c,d)$ na $(a,c)$ i $(b,d)$:
\[
c_{ab}+c_{cd} > c_{ac}+c_{bd}.
\]

\subsubsection{Aproksymacja dla metrycznego TSP}

Dla metrycznego TSP istnieją algorytmy z gwarancją:
\begin{itemize}
  \item algorytm podwajania drzewa MST daje trasę nie gorszą niż $2\cdot OPT$,
  \item algorytm Christofidesa daje gwarancję $1.5\cdot OPT$ dla symetrycznego metrycznego TSP.
\end{itemize}

Christofides: budujemy MST, znajdujemy wierzchołki nieparzystego stopnia, liczymy minimalne skojarzenie doskonałe na nich, dodajemy je do MST, otrzymujemy multigraf eulerowski, przechodzimy cyklem Eulera i skracamy powtórzenia dzięki nierówności trójkąta.

\subsubsection{Metaheurystyki}

Dla dużych instancji stosuje się algorytmy przybliżone:
\begin{itemize}
  \item symulowane wyżarzanie,
  \item tabu search,
  \item algorytmy genetyczne z reprezentacją permutacyjną,
  \item ant colony optimization,
  \item iterated local search,
  \item MCTS albo metody uczenia przez wzmacnianie w wariantach konstrukcyjnych.
\end{itemize}

Kluczowy problem reprezentacji: operator mutacji albo krzyżowania musi zachować permutację miast. Dlatego używa się np. PMX, OX, CX, swap mutation, inversion mutation.

\subsubsection{Co powiedzieć na obronie}

TSP polega na znalezieniu najkrótszego cyklu odwiedzającego każde miasto dokładnie raz. Jest NP-trudny, a wersja decyzyjna jest NP-zupełna. Dokładne metody to brute force, Held-Karp, branch and bound oraz programowanie całkowitoliczbowe. Praktycznie używa się heurystyk konstrukcyjnych, 2-opt/3-opt, Lin-Kernighana i metaheurystyk. Dla metrycznego TSP są algorytmy aproksymacyjne, np. Christofides z gwarancją $3/2$.

\subsection{Zadania interpolacji i zastosowanie interpolacji}

Interpolacja polega na wyznaczeniu funkcji, która przechodzi dokładnie przez dane punkty pomiarowe:
\[
f(x_i)=y_i,\quad i=0,\dots,n.
\]
W odróżnieniu od aproksymacji, interpolacja wymaga dokładnego dopasowania w węzłach. Jest użyteczna, gdy wartości danych traktujemy jako dokładne albo chcemy skonstruować funkcję pomocniczą do obliczeń numerycznych.

\subsubsection{Interpolacja wielomianowa}

Dla $n+1$ różnych węzłów $x_0,\dots,x_n$ istnieje dokładnie jeden wielomian $p_n$ stopnia co najwyżej $n$, taki że $p_n(x_i)=y_i$. Można go zapisać w postaci Lagrange'a:
\[
p_n(x)=\sum_{i=0}^{n}y_i L_i(x),
\quad
L_i(x)=\prod_{\substack{j=0\\j\ne i}}^{n}\frac{x-x_j}{x_i-x_j}.
\]
Bazy $L_i$ mają własność $L_i(x_j)=1$ dla $i=j$ i $0$ dla $i\ne j$.

Wzór Newtona używa ilorazów różnicowych:
\[
p_n(x)=a_0+a_1(x-x_0)+a_2(x-x_0)(x-x_1)+\dots+a_n\prod_{j=0}^{n-1}(x-x_j),
\]
gdzie $a_k=f[x_0,\dots,x_k]$. Zaletą postaci Newtona jest łatwe dodanie nowego węzła bez przeliczania całego wielomianu.

\subsubsection{Błąd interpolacji}

Jeśli funkcja $f$ ma pochodną rzędu $n+1$, to dla interpolacji wielomianowej:
\[
f(x)-p_n(x)=\frac{f^{(n+1)}(\xi)}{(n+1)!}\prod_{i=0}^{n}(x-x_i)
\]
dla pewnego $\xi$ zależnego od $x$. Błąd zależy więc od gładkości funkcji i rozmieszczenia węzłów.

\subsubsection{Efekt Rungego i dobór węzłów}

Interpolacja wielomianem wysokiego stopnia na równoodległych węzłach może silnie oscylować przy końcach przedziału. To efekt Rungego. Przeciwdziała się mu przez:
\begin{itemize}
  \item użycie węzłów Czebyszewa,
  \item interpolację kawałkami,
  \item splajny,
  \item aproksymację zamiast interpolacji, jeśli dane są zaszumione.
\end{itemize}

Węzły Czebyszewa zagęszczają się przy końcach przedziału, zmniejszając maksymalny wpływ iloczynu $\prod_i(x-x_i)$.

\subsubsection{Splajny}

Splajn to funkcja sklejana z wielomianów niskiego stopnia, zwykle trzeciego. Splajn kubiczny na przedziałach $[x_i,x_{i+1}]$:
\begin{itemize}
  \item interpoluje dane w węzłach,
  \item jest ciągły,
  \item ma ciągłą pierwszą i drugą pochodną,
  \item unika globalnych oscylacji typowych dla wielomianów wysokiego stopnia.
\end{itemize}

Warunki brzegowe mogą być naturalne ($S''(x_0)=S''(x_n)=0$), z zadanymi pochodnymi albo okresowe.

\subsubsection{Zastosowania interpolacji}

\begin{itemize}
  \item rekonstrukcja wartości między pomiarami,
  \item grafika komputerowa i animacja,
  \item przetwarzanie sygnałów i obrazów,
  \item numeryczne całkowanie i różniczkowanie,
  \item metoda elementów skończonych,
  \item tablice funkcji specjalnych,
  \item kalibracja czujników,
  \item trajektorie robotów.
\end{itemize}

\subsubsection{Co powiedzieć na obronie}

Interpolacja buduje funkcję przechodzącą przez zadane punkty. Najważniejsze metody to interpolacja wielomianowa Lagrange'a i Newtona oraz splajny. Trzeba znać istnienie i jednoznaczność wielomianu interpolacyjnego, wzór błędu oraz problem Rungego. W praktyce dla wielu punktów preferuje się splajny albo węzły Czebyszewa.

\subsection{Metody skończone rozwiązywania układów równań liniowych}

Układ równań liniowych zapisujemy jako
\[
Ax=b,
\]
gdzie $A\in\R^{n\times n}$, $b\in\R^n$, a $x\in\R^n$ jest niewiadomą. Metody skończone, nazywane też bezpośrednimi, w arytmetyce dokładnej dają rozwiązanie po skończonej liczbie operacji. Odróżnia się je od metod iteracyjnych, które generują ciąg przybliżeń.

\subsubsection{Eliminacja Gaussa}

Eliminacja Gaussa sprowadza macierz do postaci trójkątnej górnej przez zerowanie elementów pod przekątną. Następnie rozwiązuje się układ przez podstawianie wsteczne.

Dla kroku $k$ eliminujemy elementy $a_{ik}$ dla $i>k$:
\[
m_{ik}=\frac{a_{ik}}{a_{kk}},\qquad
R_i\leftarrow R_i-m_{ik}R_k.
\]
Złożoność czasowa wynosi $\Theta(n^3)$, a podstawianie wsteczne $\Theta(n^2)$.

\subsubsection{Wybór elementu głównego}

Bez pivotingu eliminacja może być niestabilna numerycznie albo niemożliwa, gdy pivot $a_{kk}=0$. W częściowym wyborze elementu głównego zamieniamy wiersze tak, aby w kolumnie $k$ wybrać największy co do modułu element wśród wierszy $k,\dots,n$. Poprawia to stabilność i w praktyce jest standardem.

Pełny pivoting pozwala zamieniać także kolumny, ale jest droższy i rzadziej potrzebny.

\subsubsection{Rozkład LU}

Eliminacja Gaussa odpowiada rozkładowi:
\[
PA=LU,
\]
gdzie $P$ jest macierzą permutacji, $L$ trójkątną dolną, a $U$ trójkątną górną. Rozwiązanie $Ax=b$ przebiega wtedy:
\[
Ly=Pb,\qquad Ux=y.
\]
Rozkład LU jest szczególnie opłacalny, gdy trzeba rozwiązać wiele układów z tą samą macierzą $A$ i różnymi prawymi stronami $b$.

\subsubsection{Rozkład Cholesky'ego}

Jeśli $A$ jest symetryczna i dodatnio określona, można użyć rozkładu:
\[
A=LL^T.
\]
Jest około dwa razy tańszy niż LU i stabilny dla tej klasy macierzy. Występuje często w metodzie najmniejszych kwadratów, procesach Gaussowskich, optymalizacji wypukłej i równaniach normalnych, choć same równania normalne mogą pogarszać uwarunkowanie.

\subsubsection{Rozkład QR}

Rozkład QR:
\[
A=QR,
\]
gdzie $Q$ jest ortogonalna, a $R$ trójkątna górna. Jest bardzo ważny w rozwiązywaniu problemów najmniejszych kwadratów:
\[
\min_x \|Ax-b\|_2.
\]
Ponieważ macierze ortogonalne nie zmieniają normy euklidesowej, QR jest numerycznie stabilniejszy niż równania normalne $A^TAx=A^Tb$.

\subsubsection{Uwarunkowanie układu}

Wrażliwość rozwiązania na zaburzenia danych opisuje liczba uwarunkowania:
\[
\kappa(A)=\|A\|\|A^{-1}\|.
\]
Jeśli $\kappa(A)$ jest duża, mały błąd w $A$ albo $b$ może spowodować duży błąd w $x$. Stabilny algorytm nie naprawi złego uwarunkowania problemu, ale nie powinien znacząco dokładać własnego błędu.

\subsubsection{Co powiedzieć na obronie}

Metody skończone rozwiązują układ po skończonej liczbie operacji. Najważniejsza jest eliminacja Gaussa z pivotingiem oraz rozkłady LU, Cholesky'ego i QR. LU jest dobre dla wielu prawych stron, Cholesky dla macierzy symetrycznych dodatnio określonych, QR dla najmniejszych kwadratów i stabilności. Trzeba wspomnieć o koszcie $\Theta(n^3)$ i uwarunkowaniu.

\subsection{Metody poszukiwania zer funkcji jednej zmiennej}

Szukamy $x^*$ takiego, że
\[
f(x^*)=0.
\]
Metody różnią się wymaganiami wobec funkcji, szybkością zbieżności i gwarancjami.

\subsubsection{Metoda bisekcji}

Zakładamy, że $f$ jest ciągła na $[a,b]$ i $f(a)f(b)<0$. Z twierdzenia Darboux istnieje zero w przedziale. W każdym kroku bierzemy środek
\[
c=\frac{a+b}{2}
\]
i wybieramy połowę przedziału, w której następuje zmiana znaku.

Zalety: prostota, gwarancja zbieżności, odporność. Wady: zbieżność liniowa i potrzeba przedziału ze zmianą znaku. Po $k$ krokach długość przedziału wynosi $(b-a)/2^k$, więc łatwo ustalić liczbę iteracji potrzebną dla tolerancji.

\subsubsection{Metoda Newtona}

Metoda Newtona korzysta z pochodnej i linearyzacji funkcji:
\[
x_{k+1}=x_k-\frac{f(x_k)}{f'(x_k)}.
\]
Interpretacja geometryczna: bierzemy styczną do wykresu w $x_k$ i jej przecięcie z osią $x$.

Dla dobrego punktu startowego i prostego pierwiastka metoda ma zbieżność kwadratową. Może jednak zawieść, jeśli $f'(x_k)$ jest bliskie zeru, punkt startowy jest zły, funkcja ma osobliwości albo iteracje opuszczają sensowny obszar.

\subsubsection{Metoda siecznych}

Metoda siecznych zastępuje pochodną ilorazem różnicowym:
\[
x_{k+1}=x_k-f(x_k)\frac{x_k-x_{k-1}}{f(x_k)-f(x_{k-1})}.
\]
Nie wymaga obliczania pochodnej. Zbieżność jest superliniowa, rzędu około $1.618$, ale brak tak silnych gwarancji jak w bisekcji.

\subsubsection{Regula falsi}

Regula falsi także używa siecznej, ale zachowuje przedział, w którym funkcja zmienia znak. Dzięki temu ma gwarancję podobną do bisekcji, choć może zbiegać wolno, jeśli jeden koniec przedziału pozostaje długo nieruchomy. Istnieją modyfikacje, np. Illinois, poprawiające ten problem.

\subsubsection{Metody hybrydowe}

W praktyce często łączy się bezpieczeństwo bisekcji z szybkością Newtona albo siecznych. Przykładem jest metoda Brenta, która używa bisekcji, siecznych i interpolacji odwrotnej kwadratowej. Jest standardowym wyborem dla jednowymiarowego znajdowania pierwiastków, gdy mamy przedział ze zmianą znaku.

\subsubsection{Kryteria stopu}

Typowe kryteria:
\begin{itemize}
  \item $|f(x_k)|<\eps$,
  \item $|x_{k+1}-x_k|<\eps$,
  \item długość przedziału $<\eps$,
  \item przekroczenie maksymalnej liczby iteracji.
\end{itemize}
Samo małe $|f(x_k)|$ może być mylące dla funkcji bardzo płaskiej albo źle przeskalowanej, dlatego często łączy się kilka kryteriów.

\subsubsection{Co powiedzieć na obronie}

Bisekcja jest wolna, ale gwarantowana przy ciągłości i zmianie znaku. Newton jest bardzo szybki lokalnie, ale wymaga pochodnej i dobrego startu. Sieczne nie wymagają pochodnej i mają zbieżność superliniową. W praktyce dobre są metody hybrydowe, np. Brent, bo łączą gwarancję z szybkością.

\subsection{Metody całkowania numerycznego}

Całkowanie numeryczne, czyli kwadratura, przybliża
\[
I=\int_a^b f(x)\,dx
\]
za pomocą skończonej liczby wartości funkcji. Stosuje się je, gdy funkcji nie da się scałkować analitycznie, znamy ją tylko z pomiarów albo całka jest składnikiem większego algorytmu.

\subsubsection{Złożone wzory Newtona-Cotesa}

Najprostsze metody dzielą przedział na podprzedziały długości $h=(b-a)/n$.

Metoda prostokątów:
\[
I\approx h\sum_{i=0}^{n-1} f(x_i)
\]
albo wariant ze środkiem:
\[
I\approx h\sum_{i=0}^{n-1} f\left(\frac{x_i+x_{i+1}}{2}\right).
\]

Metoda trapezów:
\[
I\approx h\left(\frac{f(a)+f(b)}{2}+\sum_{i=1}^{n-1}f(x_i)\right).
\]

Metoda Simpsona dla parzystego $n$:
\[
I\approx \frac{h}{3}\left(f(x_0)+f(x_n)+4\sum_{\substack{i=1\\i\text{ nieparzyste}}}^{n-1}f(x_i)+2\sum_{\substack{i=2\\i\text{ parzyste}}}^{n-2}f(x_i)\right).
\]

Dla funkcji dostatecznie gładkich złożona metoda trapezów ma błąd rzędu $\BigO(h^2)$, a metoda Simpsona $\BigO(h^4)$.

\subsubsection{Kwadratury Gaussa}

Kwadratury Gaussa dobierają węzły i wagi tak, aby całkować dokładnie wielomiany możliwie wysokiego stopnia:
\[
\int_a^b f(x)w(x)\,dx\approx \sum_{i=1}^{n} A_i f(x_i).
\]
Dla $n$ węzłów kwadratura Gaussa jest dokładna dla wielomianów stopnia do $2n-1$. Węzły są związane z pierwiastkami wielomianów ortogonalnych. Najczęstsza jest Gauss-Legendre dla wagi $w(x)=1$ na $[-1,1]$.

\subsubsection{Metody adaptacyjne}

Jeśli funkcja ma obszary trudne, np. ostre piki albo szybkie zmiany, równomierna siatka marnuje punkty tam, gdzie funkcja jest gładka. Metody adaptacyjne dzielą przedział tam, gdzie lokalny błąd jest duży. Adaptacyjny Simpson porównuje wynik na całym przedziale z sumą wyników na połowach i rekurencyjnie zagęszcza podział.

\subsubsection{Ekstrapolacja Richardsona i Romberg}

Jeśli znamy rozwinięcie błędu metody, można połączyć wyniki dla kroków $h$ i $h/2$, aby wyeliminować główny składnik błędu. Metoda Romberga stosuje ekstrapolację Richardsona do metody trapezów, tworząc tablicę coraz dokładniejszych przybliżeń.

\subsubsection{Monte Carlo}

Dla całek wielowymiarowych metody deterministyczne cierpią na przekleństwo wymiarowości. Monte Carlo przybliża wartość oczekiwaną:
\[
\int_D f(x)\,dx = |D|\mathbb{E}[f(X)],
\]
gdzie $X$ jest losowane jednostajnie z $D$. Estymator:
\[
\hat I=\frac{|D|}{N}\sum_{i=1}^{N}f(X_i)
\]
ma błąd standardowy malejący jak $\BigO(N^{-1/2})$, niezależnie od wymiaru w sensie tempa względem liczby próbek. Wadą jest wolna zbieżność i losowa zmienność.

\subsubsection{Źródła błędów}

Błąd całkowania obejmuje:
\begin{itemize}
  \item błąd obcięcia wynikający z przybliżenia całki sumą,
  \item błąd zaokrągleń arytmetyki zmiennoprzecinkowej,
  \item błędy danych wejściowych,
  \item trudności przy osobliwościach, nieciągłościach i funkcjach oscylacyjnych.
\end{itemize}

\subsubsection{Co powiedzieć na obronie}

Podstawowe metody to prostokąty, trapezy i Simpson, dalej kwadratury Gaussa, metody adaptacyjne, Romberg oraz Monte Carlo. Simpson jest dokładniejszy od trapezów dla funkcji gładkich, Gauss wykorzystuje optymalny dobór węzłów, metody adaptacyjne zagęszczają siatkę tam, gdzie funkcja jest trudna, a Monte Carlo jest szczególnie ważne w dużym wymiarze.

\section{Programowanie matematyczne i algorytmy zaawansowane}

\subsection{Metody poszukiwania ekstremum funkcji nieliniowej}

Rozważamy problem
\[
\min_{x\in\R^n} f(x)
\]
albo równoważnie maksymalizację przez minimalizację $-f(x)$. Funkcja może być wypukła albo niewypukła, gładka albo niegładka, tania albo kosztowna w ewaluacji. Wybór metody zależy głównie od dostępności pochodnych, wymiaru, wypukłości, szumu i kosztu obliczenia funkcji.

\subsubsection{Warunki optymalności bez ograniczeń}

Jeśli $f$ jest różniczkowalna i $x^*$ jest lokalnym minimum wewnętrznym, to konieczny warunek pierwszego rzędu:
\[
\nabla f(x^*)=0.
\]
Jeśli $f$ jest dwa razy różniczkowalna, konieczny warunek drugiego rzędu:
\[
\nabla^2 f(x^*) \succeq 0.
\]
Warunek wystarczający dla ścisłego minimum lokalnego:
\[
\nabla f(x^*)=0,\qquad \nabla^2 f(x^*)\succ 0.
\]
Dla funkcji wypukłej każdy punkt spełniający $\nabla f(x^*)=0$ jest minimum globalnym. Dla funkcji ściśle wypukłej minimum globalne jest co najwyżej jedno.

\subsubsection{Metoda największego spadku}

Metoda największego spadku, czyli gradient descent, wykonuje kroki:
\[
x_{k+1}=x_k-\alpha_k\nabla f(x_k),
\]
gdzie $\alpha_k>0$ jest długością kroku. Kierunek $-\nabla f(x_k)$ jest kierunkiem najszybszego lokalnego spadku w normie euklidesowej.

Długość kroku można dobrać:
\begin{itemize}
  \item jako stałą,
  \item przez dokładne przeszukiwanie liniowe,
  \item przez warunki Armijo albo Wolfe'a,
  \item według harmonogramu malejącego, zwłaszcza w metodach stochastycznych.
\end{itemize}

Dla funkcji wypukłej z gradientem Lipschitza, przy odpowiednim kroku metoda zbiega do minimum. Dla funkcji silnie wypukłej zbieżność jest liniowa, ale może być bardzo wolna przy złym uwarunkowaniu, gdy poziomice są wydłużone.

\subsubsection{Metoda Newtona}

Metoda Newtona używa przybliżenia kwadratowego:
\[
f(x_k+p)\approx f(x_k)+\nabla f(x_k)^Tp+\frac{1}{2}p^T\nabla^2f(x_k)p.
\]
Kierunek Newtona spełnia:
\[
\nabla^2 f(x_k)p_k=-\nabla f(x_k),
\]
a aktualizacja:
\[
x_{k+1}=x_k+\alpha_k p_k.
\]

Jeśli Hessian jest dodatnio określony i punkt startowy jest dostatecznie blisko minimum, metoda ma zbieżność kwadratową. Wady: koszt obliczenia i odwracania/rozwiązywania układu z Hessianem, problemy przy Hessianie osobliwym albo nieokreślonym oraz brak globalnych gwarancji bez line search albo trust region.

\subsubsection{Metody quasi-Newtonowskie}

Metody quasi-Newtonowskie przybliżają Hessian albo jego odwrotność na podstawie kolejnych gradientów. Najbardziej znana jest BFGS. Niech
\[
s_k=x_{k+1}-x_k,\qquad y_k=\nabla f(x_{k+1})-\nabla f(x_k).
\]
BFGS aktualizuje przybliżenie Hessianu tak, aby spełniało równanie secant:
\[
B_{k+1}s_k=y_k.
\]
L-BFGS przechowuje tylko kilka ostatnich par $(s_k,y_k)$, dlatego nadaje się do dużych wymiarów.

Zalety: zwykle szybsze niż gradient descent, bez jawnego Hessianu. Wady: wymagają gradientu, a dla funkcji silnie niewypukłych lub zaszumionych mogą wymagać zabezpieczeń.

\subsubsection{Metody bezpochodne}

Gdy gradient nie jest dostępny, funkcja jest szumowa albo niedokładna, stosuje się metody bezpochodne:
\begin{description}
  \item[Nelder-Mead] utrzymuje sympleks $n+1$ punktów i wykonuje odbicie, ekspansję, kontrakcję i redukcję. Działa dobrze w małych wymiarach, ale nie ma silnych gwarancji w ogólności.
  \item[Pattern search] bada skończony zbiór kierunków i zmniejsza krok, gdy nie ma poprawy.
  \item[CMA-ES] adaptuje rozkład normalny próbkowania i macierz kowariancji. Jest silny dla trudnych, niewypukłych problemów ciągłych, ale kosztowny.
  \item[Symulowane wyżarzanie] akceptuje czasem gorsze ruchy z prawdopodobieństwem malejącym wraz z temperaturą, co pomaga wyjść z minimów lokalnych.
\end{description}

\subsubsection{Optymalizacja stochastyczna}

Dla funkcji będącej średnią po danych:
\[
f(x)=\frac{1}{m}\sum_{i=1}^{m}\ell_i(x)
\]
pełny gradient może być kosztowny. Stochastic Gradient Descent używa losowej próbki albo mini-batcha:
\[
x_{k+1}=x_k-\alpha_k g_k,\qquad \mathbb{E}[g_k]=\nabla f(x_k).
\]
W uczeniu maszynowym popularne są modyfikacje z momentum, RMSProp, Adam i AdamW. Ich zaletą jest skalowalność, wadą wrażliwość na hiperparametry i trudniejsza analiza końcowej dokładności.

\subsubsection{Lokalne i globalne optimum}

Dla funkcji niewypukłych metoda lokalna zwykle gwarantuje co najwyżej znalezienie punktu stacjonarnego albo lokalnego minimum. Globalna optymalizacja wymaga dodatkowych technik: wielostartu, branch and bound, relaksacji wypukłych, metaheurystyk, metod przedziałowych albo specjalnej struktury problemu.

\subsubsection{Co powiedzieć na obronie}

Najważniejsze metody bez ograniczeń to gradient descent, Newton, quasi-Newton, metody bezpochodne i stochastyczne. Gradient descent jest prosty, ale może być wolny; Newton jest szybki lokalnie, ale kosztowny; BFGS/L-BFGS to praktyczny kompromis; Nelder-Mead i inne metody bezpochodne są użyteczne, gdy nie mamy gradientu. W problemach wypukłych optimum lokalne jest globalne, w niewypukłych nie.

\subsection{Metody poszukiwania ekstremum w obecności ograniczeń}

Problem z ograniczeniami ma postać:
\[
\min_x f(x)
\]
przy warunkach
\[
g_i(x)\le 0,\quad i=1,\dots,m,\qquad h_j(x)=0,\quad j=1,\dots,p.
\]
Zbiór dopuszczalny jest zbiorem punktów spełniających ograniczenia. Ograniczenia mogą być liniowe, nieliniowe, wypukłe, niewypukłe, równościowe, nierównościowe albo całkowitoliczbowe.

\subsubsection{Warunki KKT}

Dla problemu gładkiego definiujemy Lagrangian:
\[
\mathcal{L}(x,\lambda,\mu)=f(x)+\sum_{i=1}^{m}\lambda_i g_i(x)+\sum_{j=1}^{p}\mu_j h_j(x),
\]
gdzie $\lambda_i\ge 0$ są mnożnikami dla nierówności, a $\mu_j$ dla równości.

Warunki Karusha-Kuhna-Tuckera:
\begin{align*}
\nabla_x \mathcal{L}(x^*,\lambda^*,\mu^*)&=0 &&\text{stacjonarność},\\
g_i(x^*)&\le 0,\quad h_j(x^*)=0 &&\text{dopuszczalność pierwotna},\\
\lambda_i^*&\ge 0 &&\text{dopuszczalność dualna},\\
\lambda_i^* g_i(x^*)&=0 &&\text{komplementarność}.
\end{align*}
Dla problemów wypukłych przy odpowiednich warunkach regularności KKT są nie tylko konieczne, ale też wystarczające.

\subsubsection{Metoda mnożników Lagrange'a}

Dla samych równości $h_j(x)=0$ szukamy punktów stacjonarnych Lagrangianu:
\[
\nabla f(x^*)+\sum_{j=1}^{p}\mu_j\nabla h_j(x^*)=0,\qquad h_j(x^*)=0.
\]
Geometrycznie gradient funkcji celu w optimum leży w rozpięciu gradientów aktywnych ograniczeń, czyli nie ma składowej pozwalającej poprawić wynik wzdłuż zbioru dopuszczalnego.

\subsubsection{Metody kar i barier}

Metody kar zamieniają problem ograniczony na sekwencję problemów bez ograniczeń:
\[
\min_x f(x)+\rho_k P(x),
\]
gdzie $P(x)$ karze naruszenie ograniczeń, np.
\[
P(x)=\sum_i \max(0,g_i(x))^2+\sum_j h_j(x)^2.
\]
Gdy $\rho_k$ rośnie, rozwiązania powinny zbliżać się do dopuszczalnych. Wadą jest pogarszające się uwarunkowanie dla dużych kar.

Metody barier działają wewnątrz zbioru dopuszczalnego, np. dla $g_i(x)<0$:
\[
\min_x f(x)-\mu\sum_i \log(-g_i(x)).
\]
Gdy $\mu\to 0$, bariera słabnie, a rozwiązanie zbliża się do optimum problemu oryginalnego. To idea metod punktu wewnętrznego.

\subsubsection{Programowanie sekwencyjne kwadratowe}

Sequential Quadratic Programming rozwiązuje w każdej iteracji przybliżony problem kwadratowy: funkcję celu zastępuje modelem kwadratowym Lagrangianu, a ograniczenia linearyzuje. SQP jest bardzo skuteczne dla gładkich problemów nieliniowych średniego rozmiaru.

\subsubsection{Metody projekcyjne}

Dla prostego zbioru dopuszczalnego $C$ można wykonać krok gradientowy i projekcję:
\[
x_{k+1}=P_C(x_k-\alpha_k\nabla f(x_k)),
\]
gdzie
\[
P_C(y)=\argmin_{x\in C}\|x-y\|_2.
\]
Metoda jest naturalna dla ograniczeń pudełkowych, kul, sympleksu albo prostych zbiorów wypukłych. W uczeniu maszynowym odpowiada np. obcinaniu wag do zakresu albo projekcji na ograniczenie normy.

\subsubsection{Active set i interior point}

Metody aktywnego zbioru zgadują, które ograniczenia nierównościowe są aktywne w optimum, czyli spełnione jako równości. Następnie rozwiązują problem z tymi ograniczeniami jak z równościami i aktualizują zbiór aktywny.

Metody punktu wewnętrznego poruszają się we wnętrzu zbioru dopuszczalnego i używają barier. Są podstawą bardzo wydajnych solverów dla programowania liniowego, kwadratowego i wypukłego.

\subsubsection{Ograniczenia całkowitoliczbowe}

Jeśli część zmiennych musi być całkowita albo binarna, problem staje się programowaniem całkowitoliczbowym. Typowe techniki:
\begin{itemize}
  \item branch and bound,
  \item branch and cut,
  \item relaksacja liniowa albo wypukła,
  \item płaszczyzny tnące,
  \item heurystyki naprawcze.
\end{itemize}
Tego typu problemy są często NP-trudne.

\subsubsection{Co powiedzieć na obronie}

Dla ograniczeń kluczowe są warunki KKT: stacjonarność, dopuszczalność, nieujemność mnożników i komplementarność. Metody praktyczne to mnożniki Lagrange'a, kary, bariery, projekcja, SQP, active set i interior point. Dla problemów wypukłych warunki KKT dają globalną charakterystykę optimum; dla niewypukłych zwykle tylko lokalną.

\subsection{Wielomianowy schemat aproksymacyjny}

Wiele problemów NP-trudnych jest zbyt kosztownych do dokładnego rozwiązania dla dużych instancji. Algorytmy aproksymacyjne zwracają rozwiązania bliskie optymalnym z formalną gwarancją jakości.

\subsubsection{Współczynnik aproksymacji}

Dla problemu minimalizacyjnego algorytm ma współczynnik $\rho(n)\ge 1$, jeśli dla każdej instancji:
\[
A(I)\le \rho(n)\cdot OPT(I),
\]
gdzie $A(I)$ jest wartością rozwiązania algorytmu, a $OPT(I)$ optimum. Dla problemu maksymalizacyjnego zwykle wymaga się:
\[
A(I)\ge \frac{1}{\rho(n)}OPT(I).
\]
Jeśli $\rho$ jest stałe, mówimy o stałej aproksymacji.

\subsubsection{PTAS}

Polynomial-Time Approximation Scheme to rodzina algorytmów, która dla dowolnego $\eps>0$ zwraca rozwiązanie o jakości:
\[
A(I)\le (1+\eps)OPT(I)
\]
dla minimalizacji albo
\[
A(I)\ge (1-\eps)OPT(I)
\]
dla maksymalizacji, w czasie wielomianowym względem rozmiaru wejścia $n$ dla ustalonego $\eps$.

Istotne: czas może być bardzo zły jako funkcja $1/\eps$, np.
\[
\BigO(n^{1/\eps}).
\]
Dla każdego ustalonego $\eps$ jest to wielomian w $n$, więc formalnie PTAS.

\subsubsection{FPTAS}

Fully Polynomial-Time Approximation Scheme wymaga czasu wielomianowego zarówno względem $n$, jak i $1/\eps$, np.
\[
\BigO(n^3/\eps^2).
\]
FPTAS jest znacznie silniejszy praktycznie i teoretycznie niż PTAS.

\subsubsection{EPTAS}

Efficient PTAS ma czas postaci
\[
f(1/\eps)\cdot n^{O(1)},
\]
gdzie wykładnik przy $n$ nie zależy od $\eps$. To kompromis między PTAS i FPTAS.

\subsubsection{Przykłady}

\begin{itemize}
  \item Problem plecakowy 0/1 ma FPTAS oparty na skalowaniu wartości przedmiotów i programowaniu dynamicznym.
  \item Euklidesowy TSP w ustalonym wymiarze ma PTAS.
  \item Dla ogólnego metrycznego TSP znany jest algorytm Christofidesa z gwarancją $3/2$, ale nie jest to PTAS.
  \item Dla wielu problemów silnie NP-trudnych istnienie FPTAS implikowałoby $\mathrm{P}=\NP$.
\end{itemize}

\subsubsection{Schemat FPTAS dla plecaka}

W plecaku maksymalizujemy wartość przy ograniczeniu wagi. Klasyczne DP po wartościach działa w czasie zależnym od sumy wartości, która może być pseudowielomianowa. FPTAS skaluje wartości:
\[
v_i'=\left\lfloor \frac{v_i}{K}\right\rfloor
\]
dla odpowiednio dobranego $K$ zależnego od $\eps$ i $v_{\max}$. Następnie rozwiązujemy problem DP dla mniejszych wartości. Tracimy co najwyżej kontrolowany ułamek optimum.

\subsubsection{Co powiedzieć na obronie}

PTAS to rodzina algorytmów, która dla dowolnej dokładności $\eps$ daje rozwiązanie w granicach $1+\eps$ od optimum w czasie wielomianowym dla ustalonego $\eps$. FPTAS jest silniejszy, bo czas jest wielomianowy także względem $1/\eps$. Schematy aproksymacyjne są sposobem radzenia sobie z NP-trudnością, gdy optimum dokładne jest za drogie.

\subsection{Programowanie dynamiczne}

Programowanie dynamiczne jest techniką projektowania algorytmów dla problemów z optymalną podstrukturą i nakładającymi się podproblemami. Zamiast wielokrotnie rozwiązywać te same podproblemy, zapamiętuje się ich wyniki.

\subsubsection{Dwie podstawowe własności}

\begin{description}
  \item[Optymalna podstruktura] rozwiązanie optymalne problemu można zbudować z rozwiązań optymalnych podproblemów.
  \item[Nakładające się podproblemy] rekurencyjne rozwiązanie generowałoby te same podproblemy wielokrotnie.
\end{description}

Jeśli podproblemy się nie nakładają, zwykle wystarcza dziel i zwyciężaj, jak w mergesort. Jeśli nie ma optymalnej podstruktury, DP może dawać błędne rozwiązania.

\subsubsection{Top-down i bottom-up}

Podejście top-down z memoizacją zapisuje wyniki funkcji rekurencyjnej. Jest łatwe do napisania i oblicza tylko potrzebne stany.

Podejście bottom-up wypełnia tablicę w porządku zgodnym z zależnościami. Zwykle daje mniejszy narzut i łatwiej kontrolować pamięć.

\subsubsection{Schemat projektowania DP}

\begin{enumerate}
  \item Zdefiniować stan, czyli co opisuje podproblem.
  \item Określić wartość przechowywaną w stanie.
  \item Wyprowadzić rekurencję przejścia.
  \item Ustalić przypadki bazowe.
  \item Ustalić kolejność obliczeń.
  \item Odtworzyć rozwiązanie, jeśli potrzebna jest struktura, a nie tylko wartość.
  \item Oszacować liczbę stanów i koszt przejścia.
\end{enumerate}

\subsubsection{Przykład: najdłuższy wspólny podciąg}

Dla napisów $X[1..m]$ i $Y[1..n]$ definiujemy $dp[i][j]$ jako długość LCS prefiksów $X[1..i]$ i $Y[1..j]$:
\[
dp[i][j]=
\begin{cases}
0, & i=0\text{ lub }j=0,\\
dp[i-1][j-1]+1, & X_i=Y_j,\\
\max(dp[i-1][j],dp[i][j-1]), & X_i\ne Y_j.
\end{cases}
\]
Złożoność czasu i pamięci wynosi $\Theta(mn)$, choć pamięć dla samej długości można zredukować do $\Theta(\min(m,n))$.

\subsubsection{Przykład: plecak 0/1}

Dla przedmiotów o wagach $w_i$ i wartościach $v_i$ oraz pojemności $W$:
\[
dp[i][c]=\text{maksymalna wartość z pierwszych }i\text{ przedmiotów przy pojemności }c.
\]
Rekurencja:
\[
dp[i][c]=
\begin{cases}
dp[i-1][c], & w_i>c,\\
\max(dp[i-1][c],dp[i-1][c-w_i]+v_i), & w_i\le c.
\end{cases}
\]
Czas $\Theta(nW)$ jest pseudowielomianowy, bo zależy od wartości liczbowej $W$, a nie od liczby bitów jej zapisu.

\subsubsection{Przykład: Floyd-Warshall}

Dla najkrótszych ścieżek między wszystkimi parami:
\[
d^{(k)}_{ij}=\min\left(d^{(k-1)}_{ij},d^{(k-1)}_{ik}+d^{(k-1)}_{kj}\right).
\]
Stan oznacza najkrótszą ścieżkę z $i$ do $j$ używającą jako pośrednich tylko wierzchołków z $\{1,\dots,k\}$. Czas $\Theta(n^3)$, pamięć $\Theta(n^2)$.

\subsubsection{Przykład: Held-Karp dla TSP}

Stan:
\[
dp[S][j]=\text{koszt najkrótszej ścieżki startującej w }1,\text{ odwiedzającej }S,\text{ kończącej w }j.
\]
Rekurencja:
\[
dp[S][j]=\min_{i\in S, i\ne j} dp[S\setminus\{j\}][i]+c_{ij}.
\]
Czas $\Theta(n^2 2^n)$ jest dużo lepszy niż brute force, ale nadal wykładniczy.

\subsubsection{Pułapki}

Najczęstsze błędy:
\begin{itemize}
  \item stan nie zawiera całej informacji potrzebnej do przyszłych decyzji,
  \item rekurencja używa stanu, który nie został jeszcze obliczony,
  \item pomylenie indeksów i przypadków bazowych,
  \item nieświadome użycie pseudowielomianowej złożoności,
  \item odtworzenie rozwiązania nie jest zaplanowane,
  \item pamięć jest zbyt duża mimo akceptowalnego czasu.
\end{itemize}

\subsubsection{Co powiedzieć na obronie}

Programowanie dynamiczne stosuje się, gdy problem ma optymalną podstrukturę i nakładające się podproblemy. Definiujemy stan, rekurencję, przypadki bazowe i kolejność obliczeń. DP może być top-down z memoizacją albo bottom-up. Typowe przykłady to LCS, plecak, Floyd-Warshall, edycja napisów i Held-Karp dla TSP.

\section{Zarządzanie przedsięwzięciami informatycznymi}

\subsection{Czym charakteryzuje się projekt informatyczny?}

Projekt informatyczny jest tymczasowym przedsięwzięciem prowadzącym do wytworzenia unikalnego produktu, usługi albo rezultatu związanego z systemem informatycznym. Może to być stworzenie aplikacji, wdrożenie systemu ERP, migracja danych, integracja usług, modernizacja infrastruktury, budowa hurtowni danych albo wdrożenie modelu uczenia maszynowego.

\subsubsection{Cechy projektu}

Projekt ma:
\begin{itemize}
  \item jasno określony cel biznesowy albo techniczny,
  \item początek i koniec,
  \item ograniczony budżet, czas i zakres,
  \item interesariuszy,
  \item zespół projektowy,
  \item ryzyka i niepewność,
  \item plan realizacji,
  \item kryteria akceptacji.
\end{itemize}

W IT szczególnie istotne są niepewność wymagań, zmienność technologii, integracje z istniejącymi systemami, zależność od kompetencji ludzi oraz trudność szacowania prac twórczych.

\subsubsection{Potrójne ograniczenie i jego rozszerzenia}

Klasycznie projekt opisuje się przez trójkąt:
\[
\text{zakres} \quad-\quad \text{czas} \quad-\quad \text{koszt}.
\]
Zmiana jednego elementu wpływa na pozostałe. Jeśli zakres rośnie przy stałym terminie, trzeba zwiększyć zasoby, obniżyć jakość albo zaakceptować ryzyko opóźnienia.

W praktyce dodaje się jakość, ryzyko, satysfakcję klienta, wartość biznesową i ograniczenia organizacyjne. Projekt IT nie jest sukcesem tylko dlatego, że kod działa. Sukces oznacza dostarczenie wartości w akceptowalnym koszcie, czasie i jakości.

\subsubsection{Specyfika projektów informatycznych}

Projekty informatyczne mają kilka cech odróżniających je od wielu projektów inżynierii materialnej:
\begin{description}
  \item[Niematerialny produkt] Oprogramowania nie widać bezpośrednio, dlatego postęp bywa trudny do oceny, jeśli nie ma działających przyrostów.
  \item[Zmienność wymagań] Użytkownicy często doprecyzowują potrzeby dopiero po zobaczeniu prototypu.
  \item[Wysoka złożoność zależności] Integracje, dane, infrastruktura, bezpieczeństwo i wydajność tworzą zależności ukryte.
  \item[Niski koszt kopiowania, wysoki koszt wytworzenia] Samo powielenie programu jest tanie, ale jego zaprojektowanie, testowanie i utrzymanie jest kosztowne.
  \item[Dług techniczny] Decyzje przyspieszające krótkoterminowo mogą zwiększać koszt dalszego rozwoju.
  \item[Trudność szacowania] Zadania programistyczne są często poznawczo nowe, więc estymacje mają dużą niepewność.
\end{description}

\subsubsection{Interesariusze}

Interesariusz to osoba albo organizacja wpływająca na projekt albo dotknięta jego wynikiem. W projekcie IT są to zwykle:
\begin{itemize}
  \item sponsor biznesowy,
  \item product owner albo właściciel produktu,
  \item użytkownicy końcowi,
  \item zespół projektowy,
  \item administratorzy i zespoły utrzymaniowe,
  \item dział bezpieczeństwa,
  \item compliance i prawnicy,
  \item dostawcy zewnętrzni,
  \item klienci.
\end{itemize}

Dobre zarządzanie projektem wymaga rozumienia, kto podejmuje decyzje, kto akceptuje rezultat, kto dostarcza wymagania i kto ponosi skutki wdrożenia.

\subsubsection{Zakres, wymagania i akceptacja}

Zakres określa, co projekt ma dostarczyć. Wymagania mogą być:
\begin{itemize}
  \item funkcjonalne - co system ma robić,
  \item niefunkcjonalne - wydajność, bezpieczeństwo, dostępność, użyteczność,
  \item biznesowe - jaki problem organizacji ma zostać rozwiązany,
  \item techniczne - standardy, technologie, integracje,
  \item regulacyjne - zgodność z prawem i politykami.
\end{itemize}

Kryteria akceptacji muszą być weryfikowalne. Zamiast "system ma być szybki" lepiej zapisać "95 percent odpowiedzi API dla operacji X ma czas poniżej 300 ms przy obciążeniu Y".

\subsubsection{Ryzyka typowe dla projektów IT}

\begin{itemize}
  \item niedoszacowanie złożoności,
  \item niestabilne wymagania,
  \item brak dostępności kluczowych osób,
  \item opóźnienia integracji,
  \item niska jakość danych,
  \item problemy z bezpieczeństwem,
  \item vendor lock-in,
  \item brak środowisk testowych,
  \item niedostateczne testy,
  \item brak planu migracji i rollbacku,
  \item konflikt między terminem a jakością.
\end{itemize}

\subsubsection{Co powiedzieć na obronie}

Projekt informatyczny to tymczasowe przedsięwzięcie tworzące unikalny rezultat IT. Charakteryzuje się określonym celem, zakresem, terminem, budżetem, interesariuszami i ryzykiem. Szczególne dla IT są zmienne wymagania, niematerialność produktu, trudność szacowania, złożoność integracji i dług techniczny.

\subsection{Główne różnice między projektem a pracą operacyjną}

Projekt i praca operacyjna są potrzebne organizacji, ale mają różny charakter. Projekt wprowadza zmianę; operacje utrzymują ciągłe działanie.

\subsubsection{Porównanie}

\begin{longtable}{p{0.25\textwidth}p{0.33\textwidth}p{0.33\textwidth}}
\toprule
Kryterium & Projekt & Praca operacyjna \\
\midrule
Czas trwania & tymczasowy, ma początek i koniec & ciągła, powtarzalna \\
Cel & wytworzenie unikalnego rezultatu albo zmiany & utrzymanie działania i świadczenie usług \\
Zakres & zdefiniowany dla konkretnego przedsięwzięcia & stabilny, procesowy \\
Niepewność & zwykle wysoka & zwykle niższa, procedury są znane \\
Sukces & dostarczenie wartości w czasie, budżecie i jakości & stabilność, efektywność, SLA, jakość obsługi \\
Zespół & często czasowy i interdyscyplinarny & stały albo rotacyjny \\
Przykład IT & budowa nowej aplikacji & utrzymanie aplikacji produkcyjnej \\
\bottomrule
\end{longtable}

\subsubsection{Relacja między projektem i operacjami}

Projekt często przekazuje wynik do operacji. Przykład: zespół projektowy wdraża system, a potem zespół utrzymaniowy monitoruje go, obsługuje incydenty, wykonuje aktualizacje i wspiera użytkowników.

Granica nie zawsze jest ostra. Rozwój produktu cyfrowego może być ciągłym strumieniem zmian, ale poszczególne większe inicjatywy nadal mają charakter projektowy. W podejściu produktowym mówi się częściej o ciągłym rozwoju produktu niż o jednorazowych projektach, jednak nadal występują ograniczone inicjatywy, budżety i cele.

\subsubsection{Przykład}

Wdrożenie modułu płatności online jest projektem: ma zakres, termin, integracje, testy bezpieczeństwa i odbiór. Codzienne monitorowanie transakcji, reagowanie na błędy płatności i aktualizacja certyfikatów to operacje.

\subsubsection{Co powiedzieć na obronie}

Projekt jest tymczasowy i tworzy unikalną zmianę, a praca operacyjna jest ciągła i powtarzalna. Projekt kończy się przekazaniem rezultatu, natomiast operacje utrzymują usługę, proces albo system. W IT projekty często budują lub zmieniają systemy, a operacje zapewniają ich stabilne działanie.

\subsection{Metodyki zarządzania projektami: charakterystyka, różnice i cechy wspólne}

Metodyka zarządzania projektem określa sposób planowania, organizacji, realizacji, kontroli i zamykania prac. W IT najczęściej porównuje się podejścia predykcyjne, zwinne i hybrydowe.

\subsubsection{Waterfall}

Waterfall, czyli model kaskadowy, zakłada sekwencyjne fazy:
\begin{enumerate}
  \item analiza wymagań,
  \item projekt,
  \item implementacja,
  \item testowanie,
  \item wdrożenie,
  \item utrzymanie.
\end{enumerate}

Zalety:
\begin{itemize}
  \item jasny plan i dokumentacja,
  \item dobre dopasowanie do stałych wymagań,
  \item łatwiejsze kontraktowanie zakresu,
  \item formalna kontrola zmian.
\end{itemize}

Wady:
\begin{itemize}
  \item późna informacja zwrotna od użytkownika,
  \item kosztowne zmiany wymagań,
  \item ryzyko odkrycia problemów dopiero pod koniec,
  \item słabe dopasowanie do niepewnych produktów cyfrowych.
\end{itemize}

Waterfall jest sensowny, gdy wymagania są stabilne, środowisko regulowane, a koszt zmian wysoki, np. wybrane systemy krytyczne albo kontrakty publiczne.

\subsubsection{PMBOK i podejście procesowe}

PMBOK nie jest jedną sekwencją prac, lecz zbiorem dobrych praktyk zarządzania projektami. Obejmuje obszary takie jak zakres, harmonogram, koszt, jakość, zasoby, komunikacja, ryzyko, zamówienia i interesariusze. Jest metodycznie neutralny: można go stosować w projektach kaskadowych, zwinnych i hybrydowych.

Mocną stroną PMBOK jest kompletność zarządcza. Słabością w małych zespołach może być nadmiar formalizacji, jeśli stosuje się go bez dostosowania.

\subsubsection{PRINCE2}

PRINCE2 jest metodyką silnie nastawioną na uzasadnienie biznesowe, role, produkty zarządcze i kontrolę etapową. Kluczowe elementy:
\begin{itemize}
  \item ciągła zasadność biznesowa,
  \item uczenie się z doświadczeń,
  \item zdefiniowane role i odpowiedzialności,
  \item zarządzanie etapami,
  \item zarządzanie przez wyjątki,
  \item koncentracja na produktach,
  \item dostosowanie do środowiska projektu.
\end{itemize}

PRINCE2 dobrze pasuje do organizacji wymagających kontroli i raportowania, ale może być zbyt formalny dla małych, szybko zmieniających się produktów.

\subsubsection{Agile}

Agile to rodzina podejść opartych na iteracyjności, przyrostowym dostarczaniu wartości, współpracy z klientem i reagowaniu na zmianę. Nie oznacza braku planu. Oznacza planowanie adaptacyjne i częstą weryfikację założeń.

Wspólne idee:
\begin{itemize}
  \item krótkie iteracje,
  \item częste dostarczanie działającego przyrostu,
  \item bliska współpraca z użytkownikiem lub właścicielem produktu,
  \item priorytetyzowany backlog,
  \item retrospekcje i ciągłe usprawnianie,
  \item akceptacja zmienności wymagań.
\end{itemize}

\subsubsection{Scrum}

Scrum jest ramą pracy zwinnej. Role:
\begin{description}
  \item[Product Owner] odpowiada za wartość produktu i backlog.
  \item[Scrum Master] dba o proces, usuwa przeszkody i wspiera zespół.
  \item[Developers] wytwarzają przyrost produktu.
\end{description}

Zdarzenia:
\begin{itemize}
  \item Sprint,
  \item Sprint Planning,
  \item Daily Scrum,
  \item Sprint Review,
  \item Sprint Retrospective.
\end{itemize}

Artefakty:
\begin{itemize}
  \item Product Backlog,
  \item Sprint Backlog,
  \item Increment.
\end{itemize}

Scrum dobrze działa, gdy produkt można rozwijać przyrostowo, a wymagania są zmienne. Nie rozwiązuje sam problemów technicznych: wymaga praktyk inżynierskich, testów, CI/CD i dbania o architekturę.

\subsubsection{Kanban}

Kanban koncentruje się na przepływie pracy. Typowe zasady:
\begin{itemize}
  \item wizualizacja pracy na tablicy,
  \item limity WIP, czyli ograniczenie pracy w toku,
  \item zarządzanie przepływem,
  \item jawne polityki procesu,
  \item ciągłe doskonalenie.
\end{itemize}

Kanban jest dobry dla utrzymania, zespołów operacyjnych, obsługi zgłoszeń i środowisk, gdzie praca napływa nieregularnie. W przeciwieństwie do Scruma nie wymaga sprintów.

\subsubsection{Extreme Programming}

XP jest metodyką zwinną silnie skupioną na praktykach inżynierskich:
\begin{itemize}
  \item pair programming,
  \item test-driven development,
  \item continuous integration,
  \item refaktoryzacja,
  \item prosta architektura,
  \item małe wydania,
  \item wspólna własność kodu.
\end{itemize}

XP odpowiada na lukę często widoczną w Scrumie: sama organizacja iteracji nie gwarantuje jakości technicznej.

\subsubsection{Lean}

Lean skupia się na maksymalizacji wartości i eliminacji marnotrawstwa. W IT marnotrawstwem może być nadmiar funkcji, oczekiwanie, przełączanie kontekstu, częściowo wykonana praca, defekty, ręczne czynności i nadmierna biurokracja. Lean podkreśla optymalizację całego strumienia wartości, nie tylko lokalną produktywność osób.

\subsubsection{DevOps}

DevOps nie jest klasyczną metodyką projektową, ale zestawem praktyk łączących rozwój, utrzymanie i automatyzację dostarczania. Ważne elementy:
\begin{itemize}
  \item CI/CD,
  \item infrastructure as code,
  \item monitoring i observability,
  \item automatyzacja testów,
  \item szybki rollback,
  \item współodpowiedzialność za produkcję.
\end{itemize}

DevOps skraca dystans między projektem i operacjami.

\subsubsection{Podejścia hybrydowe}

W praktyce często łączy się elementy. Przykład: formalne etapy i budżet według PRINCE2, a realizacja zespołu w Scrumie. Albo roadmapa produktu i kwartalne cele, ale codzienny przepływ w Kanbanie. Hybryda jest sensowna, jeśli wynika z potrzeb organizacji, a nie z niekonsekwencji.

\subsubsection{Cechy wspólne metodyk}

Mimo różnic większość metodyk obejmuje:
\begin{itemize}
  \item definiowanie celu i zakresu,
  \item role i odpowiedzialności,
  \item planowanie pracy,
  \item monitorowanie postępu,
  \item zarządzanie ryzykiem,
  \item komunikację z interesariuszami,
  \item kontrolę jakości,
  \item obsługę zmian,
  \item zamknięcie albo przekazanie rezultatu.
\end{itemize}

\subsubsection{Najważniejsze różnice}

Waterfall zakłada stabilność wymagań i sekwencyjność. Agile zakłada zmienność i iteracyjne dostarczanie wartości. Scrum daje ramę pracy zespołowej w iteracjach. Kanban optymalizuje przepływ i limity pracy w toku. PRINCE2 i PMBOK są bardziej zarządcze i formalne. XP skupia się na jakości inżynierskiej. DevOps skupia się na automatyzacji dostarczania i odpowiedzialności za działanie systemu.

\subsubsection{Co powiedzieć na obronie}

Metodyki różnią się podejściem do zmiany, planowania i kontroli. Kaskadowe są dobre przy stabilnych wymaganiach, zwinne przy niepewności i potrzebie częstej informacji zwrotnej, Kanban przy ciągłym przepływie zadań, PRINCE2/PMBOK przy formalnym zarządzaniu, XP przy nacisku na praktyki programistyczne. Wspólne są cel, role, planowanie, monitorowanie, ryzyko, jakość i komunikacja.

\section{Reprezentacja wiedzy}

\subsection{Metoda rezolucji w rachunku predykatów}

Rezolucja jest regułą wnioskowania i zarazem podstawą automatycznego dowodzenia twierdzeń. Jej główna idea polega na sprowadzeniu dowodu $T\models \beta$ do wykazania sprzeczności zbioru $T\cup\{\neg\beta\}$. Jeśli założenia razem z negacją tezy są niespełnialne, to teza wynika z założeń.

\subsubsection{Idea dowodu przez sprzeczność}

Dla teorii $T$ i zdania $\beta$ zachodzi:
\[
T\models \beta \quad\Longleftrightarrow\quad T\cup\{\neg\beta\}\text{ jest niespełnialny}.
\]
Rezolucja próbuje więc wyprowadzić klauzulę pustą, oznaczaną często $\square$. Klauzula pusta reprezentuje fałsz, czyli sprzeczność.

W rachunku zdań metoda jest decyzyjna: jeśli zbiór klauzul jest niespełnialny, rezolucja może znaleźć sprzeczność; jeśli jest spełnialny, pełne przeszukanie może to wykazać. W rachunku predykatów pierwszego rzędu sytuacja jest trudniejsza: logika pierwszego rzędu jest półrozstrzygalna. Jeśli wniosek rzeczywiście wynika, procedura może znaleźć dowód; jeśli nie wynika, przeszukiwanie może nie zakończyć się.

\subsubsection{Literały i klauzule}

Atom to formuła postaci $P(t_1,\dots,t_k)$, gdzie $P$ jest predykatem, a $t_i$ są termami. Literał to atom albo negacja atomu. Klauzula jest alternatywą literałów:
\[
L_1\vee L_2\vee\dots\vee L_m.
\]
Zbiór klauzul interpretuje się jako koniunkcję wszystkich klauzul. Klauzula pusta $\square$ jest alternatywą zera literałów, czyli jest zawsze fałszywa.

\subsubsection{Reguła rezolucji w rachunku zdań}

Jeśli mamy dwie klauzule zawierające literały komplementarne:
\[
\lambda\vee A,\qquad \neg\lambda\vee B,
\]
to możemy wyprowadzić rezolwentę:
\[
A\vee B.
\]
Reguła jest poprawna semantycznie, bo jeśli pierwsza i druga klauzula są prawdziwe, to po usunięciu pary komplementarnej pozostała alternatywa musi być prawdziwa.

Przykład:
\[
P\vee Q,\qquad \neg P\vee R
\]
dają:
\[
Q\vee R.
\]

\subsubsection{Sprowadzenie formuł do postaci klauzulowej}

Aby stosować rezolucję, formuły trzeba zamienić na postać normalną. Dla rachunku zdań używa się CNF:
\[
\bigwedge_i \bigvee_j L_{ij}.
\]
Algorytm:
\begin{enumerate}
  \item usunąć implikacje i równoważności:
  \[
  \alpha\to\beta \equiv \neg\alpha\vee\beta,
  \]
  \[
  \alpha\leftrightarrow\beta \equiv (\neg\alpha\vee\beta)\wedge(\neg\beta\vee\alpha);
  \]
  \item przesunąć negacje do atomów, używając praw de Morgana;
  \item zastosować rozdzielność $\vee$ względem $\wedge$;
  \item potraktować koniunkcję jako zbiór klauzul.
\end{enumerate}

\subsubsection{Dodatkowe problemy w rachunku predykatów}

W rachunku predykatów pojawiają się kwantyfikatory i zmienne. Dwa literały mogą nie być identyczne syntaktycznie, ale mogą stać się identyczne po podstawieniu. Przykład:
\[
Student(x),\qquad Student(Piotr).
\]
Nie są tym samym literałem, ale podstawienie $\{x/Piotr\}$ je uzgadnia. Do rezolucji pierwszego rzędu potrzebujemy więc:
\begin{itemize}
  \item postaci preneksowej i klauzulowej,
  \item standaryzacji nazw zmiennych,
  \item skolemizacji,
  \item unifikacji,
  \item reguły rezolucji z najogólniejszym unifikatorem.
\end{itemize}

\subsubsection{Postać preneksowa}

Formuła jest w postaci preneksowej, gdy wszystkie kwantyfikatory stoją z przodu:
\[
Q_1x_1Q_2x_2\dots Q_nx_n.\; M,
\]
gdzie $M$ jest formułą bez kwantyfikatorów, nazywaną macierzą.

Przekształcanie do postaci preneksowej obejmuje:
\begin{enumerate}
  \item usunięcie implikacji i równoważności,
  \item przesunięcie negacji do atomów:
  \[
  \neg\forall x\,\alpha \equiv \exists x\,\neg\alpha,\qquad
  \neg\exists x\,\alpha \equiv \forall x\,\neg\alpha,
  \]
  \item zmianę nazw zmiennych związanych, aby uniknąć kolizji,
  \item wyniesienie kwantyfikatorów na początek,
  \item przekształcenie macierzy do CNF.
\end{enumerate}

\subsubsection{Skolemizacja}

Skolemizacja usuwa kwantyfikatory egzystencjalne, zastępując zmienne egzystencjalne stałymi albo funkcjami Skolema. Nie zachowuje równoważności logicznej wprost, ale zachowuje spełnialność, co wystarcza dla rezolucji refutacyjnej.

Jeśli mamy:
\[
\exists y\, P(y),
\]
zastępujemy $y$ nową stałą $c$:
\[
P(c).
\]

Jeśli zmienna egzystencjalna zależy od wcześniejszych uniwersalnych:
\[
\forall x\exists y\, Loves(x,y),
\]
to $y$ zastępujemy funkcją Skolema $f(x)$:
\[
\forall x\, Loves(x,f(x)).
\]
Po skolemizacji wszystkie pozostałe zmienne są traktowane jako uniwersalnie kwantyfikowane, więc kwantyfikatory uniwersalne można pominąć w zapisie klauzul.

\subsubsection{Unifikacja}

Unifikacja znajduje podstawienie, które czyni dwa termy albo literały identycznymi. Najważniejszy jest najogólniejszy unifikator, MGU, czyli taki unifikator, z którego każdy inny unifikator może być otrzymany przez dalsze podstawienie.

Przykłady:
\[
P(x,a)\quad\text{i}\quad P(f(y),a)
\]
mają MGU:
\[
\{x/f(y)\}.
\]

Literały:
\[
Parent(x,John),\qquad Parent(Mary,y)
\]
mają MGU:
\[
\{x/Mary,\;y/John\}.
\]

W unifikacji trzeba stosować occurs check: zmienna nie może zostać podstawiona termem zawierającym tę samą zmienną, np. $x=f(x)$ nie powinno być dopuszczone w standardowej logice pierwszego rzędu, bo tworzyłoby nieskończony term.

\subsubsection{Reguła rezolucji pierwszego rzędu}

Jeśli mamy klauzule:
\[
L\vee A,\qquad \neg K\vee B,
\]
a $\theta$ jest MGU literałów $L$ i $K$, to rezolwentą jest:
\[
(A\vee B)\theta.
\]
Najpierw unifikujemy literały komplementarne, a potem usuwamy je i stosujemy podstawienie do reszty klauzuli.

Przykład:
\[
Human(x)\to Mortal(x),\qquad Human(Socrates).
\]
Po CNF:
\[
\neg Human(x)\vee Mortal(x),\qquad Human(Socrates).
\]
Unifikujemy $Human(x)$ z $Human(Socrates)$ przez $\{x/Socrates\}$ i otrzymujemy:
\[
Mortal(Socrates).
\]
Jeśli chcemy udowodnić $Mortal(Socrates)$ przez refutację, dodajemy $\neg Mortal(Socrates)$, a potem rezolucja daje klauzulę pustą.

\subsubsection{Procedura dowodzenia rezolucyjnego}

\begin{enumerate}
  \item Weź teorię $T$ i tezę $\beta$.
  \item Dodaj negację tezy: $T\cup\{\neg\beta\}$.
  \item Przekształć wszystkie formuły do postaci klauzulowej:
  usunięcie implikacji, przesunięcie negacji, standaryzacja zmiennych, preneks, skolemizacja, CNF, usunięcie kwantyfikatorów uniwersalnych.
  \item Wybieraj pary klauzul z unifikowalnymi literałami komplementarnymi.
  \item Dodawaj rezolwenty do zbioru klauzul.
  \item Jeśli pojawi się $\square$, dowód zakończony sukcesem.
\end{enumerate}

\subsubsection{Poprawność, pełność i problem strategii}

Rezolucja jest poprawna: każda wyprowadzona klauzula wynika logicznie z poprzednich. Jest też pełna refutacyjnie dla logiki pierwszego rzędu: jeśli zbiór klauzul jest niespełnialny, istnieje dowód rezolucyjny prowadzący do klauzuli pustej.

Pełność nie oznacza efektywności. Naiwne generowanie wszystkich rezolwent prowadzi do eksplozji kombinatorycznej. Dlatego stosuje się strategie:
\begin{itemize}
  \item rezolucję liniową,
  \item set of support,
  \item preferencję klauzul jednostkowych,
  \item subsumpcję,
  \item porządkowanie literałów,
  \item indeksowanie termów,
  \item ograniczenia głębokości albo heurystyki wyboru.
\end{itemize}

\subsubsection{Rezolucja SLD i Prolog}

W programowaniu logicznym, szczególnie w Prologu, używa się wyspecjalizowanej rezolucji SLD dla klauzul Horna. Klauzula Horna ma co najwyżej jeden literał pozytywny. Reguła:
\[
Head \leftarrow Body_1,\dots,Body_k
\]
odpowiada formule:
\[
Body_1\wedge\dots\wedge Body_k\to Head.
\]
Prolog rozwiązuje zapytanie przez próbę refutacji jego negacji, wybierając cele od lewej do prawej i klauzule w kolejności programu. To jest efektywne i praktyczne, ale strategia Prologa nie jest pełna dla wszystkich możliwych programów z powodu deterministycznego porządku przeszukiwania i możliwości zapętlenia.

\subsubsection{Co powiedzieć na obronie}

Rezolucja w rachunku predykatów polega na dowodzeniu przez sprzeczność: do teorii dodaje się negację tezy, przekształca wszystko do zbioru klauzul, usuwa kwantyfikatory egzystencjalne przez skolemizację, a literały dopasowuje przez unifikację. Reguła rezolucji usuwa parę literałów komplementarnych po zastosowaniu MGU. Wyprowadzenie klauzuli pustej oznacza, że teza wynika z teorii. Metoda jest poprawna i pełna refutacyjnie, ale wymaga strategii, bo liczba rezolwent może eksplodować.

\subsection{Główne problemy w zagadnieniach wnioskowania o działaniach}

Wnioskowanie o działaniach dotyczy systemów dynamicznych: mamy obiekty, ich własności oraz akcje zmieniające stan świata. Celem jest formalne odpowiadanie na pytania typu: co będzie prawdziwe po wykonaniu akcji, czy akcja jest wykonalna, czy jakiś cel da się osiągnąć sekwencją akcji, jakie skutki są konieczne, a jakie możliwe.

\subsubsection{Model dynamicznego świata}

Stan świata opisuje wartości fluentów, czyli własności mogących zmieniać się w czasie. Akcje przekształcają stany. W najprostszym modelu mamy:
\begin{itemize}
  \item zbiór fluentów $F$,
  \item zbiór akcji $Ac$,
  \item stany jako interpretacje fluentów,
  \item stan początkowy,
  \item funkcję przejścia $Res(A,\sigma)$ zwracającą możliwe stany po wykonaniu akcji $A$ w stanie $\sigma$.
\end{itemize}

Jeśli akcja jest deterministyczna, $Res(A,\sigma)$ ma jeden element. Jeśli niedeterministyczna, może mieć wiele wyników. Jeśli akcja jest niewykonalna, wynik może być pusty albo zgodnie z przyjętą semantyką może oznaczać brak efektu.

\subsubsection{Ontologia działań}

W modelowaniu działań trzeba zdecydować, jakie typy zjawisk dopuszczamy:
\begin{itemize}
  \item skutki deterministyczne albo niedeterministyczne,
  \item skutki pewne albo typowe, czyli domyślne,
  \item skutki bezpośrednie i pośrednie,
  \item skutki natychmiastowe i opóźnione,
  \item akcje jednoetapowe albo trwające w czasie,
  \item akcje sekwencyjne albo współbieżne,
  \item akcje elementarne albo złożone,
  \item pełną albo niepełną wiedzę o stanie i akcjach.
\end{itemize}

Każda decyzja upraszcza albo komplikuje semantykę. Prosty język akcji może wystarczyć dla sekwencyjnych działań deterministycznych, ale zawiedzie przy skutkach ubocznych, współbieżności albo wyjątkach.

\subsubsection{Frame problem}

Frame problem to problem reprezentowania i wnioskowania o tym, co nie zmienia się po wykonaniu akcji. W świecie z wieloma fluentami większość własności zwykle pozostaje bez zmian. Naiwne dopisywanie dla każdej akcji i każdego fluentu aksjomatu "ta akcja nie zmienia tej własności" prowadzi do ogromnej liczby aksjomatów.

Przykład: w Yale Shooting Problem załadowanie broni zmienia fluent $loaded$, ale nie powinno zmieniać faktu, że Fred żyje. Logika klasyczna bez dodatkowej zasady inercji nie pozwala automatycznie wywnioskować, że Fred nadal żyje po załadowaniu broni.

Rozwiązaniem jest zasada inercji: fluent zachowuje wartość, dopóki akcja nie wymusza zmiany albo nie zwalnia fluentu spod inercji. Formalizacje akcji często budują tę zasadę w semantykę, zamiast wypisywać wszystkie frame axioms.

\subsubsection{Ramification problem}

Ramification problem dotyczy skutków pośrednich. Akcja może bezpośrednio zmieniać jeden fluent, ale przez ograniczenia stanu wymuszać zmiany innych fluentów.

Przykład: jeśli obowiązuje reguła $walking\to alive$, a strzał powoduje $\neg alive$, to pośrednio powinno wynikać $\neg walking$. Nie chcemy ręcznie wypisywać każdego skutku pośredniego każdej akcji.

Trudność polega na tym, aby:
\begin{itemize}
  \item reprezentować ograniczenia statyczne,
  \item propagować skutki bezpośrednie przez te ograniczenia,
  \item nie generować nieintuicyjnych skutków,
  \item odróżnić skutki pośrednie od warunków wykonalności.
\end{itemize}

W slajdach pojawia się przykład, w którym próba zachowania spójności ograniczeń może prowadzić do dziwnych wniosków, np. "ożywienia" martwego indyka przy akcji zachęcania do chodzenia. To pokazuje, że sama minimalizacja zmian bywa niewystarczająca.

\subsubsection{Qualification problem}

Qualification problem dotyczy warunków, pod którymi akcja jest możliwa i przynosi oczekiwane skutki. W praktyce liczba potencjalnych warunków jest ogromna i nie da się ich wszystkich wypisać.

Klasyczny przykład: aby uruchomić samochód, trzeba mieć kluczyk, paliwo, sprawny akumulator, sprawne przewody i bardzo wiele innych rzeczy, w tym brak ziemniaka w rurze wydechowej. Nie da się w każdym modelu jawnie sprawdzać wszystkich nietypowych przeszkód.

Rozwiązaniem mogą być:
\begin{itemize}
  \item domyślne założenia normalności,
  \item logiki niemonotoniczne,
  \item preconditions z wyjątkami,
  \item probabilistyczne modele skuteczności działań,
  \item rozróżnienie skutków typowych i nietypowych.
\end{itemize}

\subsubsection{Problem skutków domyślnych}

W realnym świecie działania mają często skutki typowe, ale niepewne. Strzał z załadowanej broni typowo zabija, ale mogą istnieć wyjątki: broń się zacina, cel ma osłonę, strzał chybia. Logika klasyczna jest monotoniczna: dodanie nowej informacji nie usuwa wcześniejszych wniosków. Rozumowanie domyślne jest niemonotoniczne: wniosek "typowo $P$" może zostać wycofany po poznaniu wyjątku.

Modele działań z efektami domyślnymi muszą odróżniać:
\begin{itemize}
  \item skutki pewne,
  \item skutki normalne albo preferowane,
  \item skutki atypowe,
  \item minimalizację odstępstw od normalności.
\end{itemize}

\subsubsection{Niedeterminizm}

Akcja może mieć wiele możliwych rezultatów. Przykład: rzut kostką, próba otwarcia uszkodzonych drzwi, ruch robota po śliskiej powierzchni. Wtedy pytania dzielą się na:
\begin{itemize}
  \item czy cel jest konieczny po akcji, czyli zachodzi we wszystkich możliwych wynikach,
  \item czy cel jest możliwy, czyli zachodzi w co najmniej jednym wyniku,
  \item czy istnieje plan gwarantujący cel,
  \item czy istnieje strategia reagująca na obserwacje.
\end{itemize}

Niedeterminizm zwiększa złożoność planowania, bo pojedyncza sekwencja akcji może nie wystarczyć; potrzebna może być strategia warunkowa.

\subsubsection{Współbieżność}

Wiele akcji może zachodzić jednocześnie. Pojawiają się pytania:
\begin{itemize}
  \item kiedy akcje są konfliktowe,
  \item czy skutki akcji złożonej są sumą skutków akcji elementarnych,
  \item czy jedna akcja blokuje drugą,
  \item czy akcje razem mają skutek inny niż osobno.
\end{itemize}

Przykład: dwie niezależne akcje, jak otwarcie drzwi i podniesienie miski, mogą mieć łączny skutek będący koniunkcją skutków. Ale dwie akcje podnoszenia miski jedną ręką mogą osobno rozlać zupę, a razem już nie. To łamie prostą zasadę dziedziczenia skutków.

\subsubsection{Czas, akcje trwające i scenariusze}

Jeśli akcje trwają w czasie, trzeba reprezentować:
\begin{itemize}
  \item moment rozpoczęcia i zakończenia,
  \item wartości fluentów w trakcie akcji,
  \item skutki opóźnione,
  \item akcje wyzwalane przez stany,
  \item akcje wyzwalane przez inne akcje,
  \item scenariusze jako zbiory wystąpień akcji w czasie.
\end{itemize}

W modelach z czasem liniowym stan w chwili $t+1$ zależy od stanu w chwili $t$, wykonanych akcji i ewentualnych reguł dynamicznych. W modelach z czasem rozgałęzionym analizuje się wiele możliwych przyszłości.

\subsubsection{Niepełna wiedza i obserwacje}

Agent często nie zna pełnego stanu świata. Wtedy zamiast pojedynczego stanu ma zbiór stanów możliwych albo przekonanie probabilistyczne. Akcje mogą być:
\begin{itemize}
  \item ontic, czyli zmieniać świat,
  \item sensing, czyli dostarczać informacji,
  \item announcement, czyli zmieniać wiedzę agentów.
\end{itemize}

Wnioskowanie musi wtedy rozróżniać prawdę w świecie od wiedzy albo przekonań agenta. Stąd powiązanie z logikami epistemicznymi, modalnymi i dynamicznymi.

\subsubsection{Planowanie}

Wnioskowanie o działaniach jest ściśle związane z planowaniem. Typowe zapytania:
\begin{itemize}
  \item Czy po sekwencji akcji $\langle A_1,\dots,A_n\rangle$ zachodzi $\alpha$?
  \item Czy istnieje sekwencja akcji prowadząca do celu?
  \item Czy każda realizacja akcji prowadzi do celu?
  \item Czy akcja jest wykonalna w danym stanie?
  \item Jakie fluenty są obserwowalne po akcji?
\end{itemize}

W modelu niedeterministycznym albo niepełnoinformacyjnym plan może być warunkowy: następna akcja zależy od obserwacji.

\subsubsection{Co powiedzieć na obronie}

Główne problemy wnioskowania o działaniach to frame problem, ramification problem i qualification problem. Frame problem dotyczy tego, co pozostaje bez zmian; ramification problem skutków pośrednich; qualification problem ogromnej liczby warunków potrzebnych do skutecznego wykonania akcji. Dodatkowo ważne są niedeterminizm, skutki domyślne, współbieżność, czas trwania akcji, niepełna wiedza i planowanie. Formalne języki akcji próbują tak dobrać semantykę, aby uzyskać intuicyjne wnioski bez ręcznego wypisywania wszystkich wyjątków i skutków.

\section{Metody sztucznej inteligencji i systemy ekspertowe}

\subsection{Sposób działania algorytmu ewolucyjnego}

Algorytm ewolucyjny jest populacyjną metaheurystyką inspirowaną ewolucją biologiczną. Przetwarza zbiór kandydatów na rozwiązania, stosuje selekcję, rekombinację i mutację, a jakość kandydatów ocenia funkcją dopasowania. Celem jest iteracyjne przesuwanie populacji w kierunku lepszych rozwiązań przy zachowaniu wystarczającej różnorodności.

\subsubsection{Elementy algorytmu}

\begin{description}
  \item[Osobnik] kandydat na rozwiązanie problemu.
  \item[Genotyp] reprezentacja osobnika używana przez operatory, np. bitstring, wektor liczb rzeczywistych, permutacja.
  \item[Fenotyp] rozwiązanie po dekodowaniu genotypu, np. konkretna trasa TSP.
  \item[Populacja] zbiór osobników.
  \item[Funkcja dopasowania] miara jakości osobnika. Dla problemu minimalizacji często przekształca się koszt na dopasowanie albo stosuje selekcję opartą bezpośrednio na rankingu.
  \item[Selekcja] wybór osobników do reprodukcji.
  \item[Krzyżowanie] tworzenie potomków z informacji wielu rodziców.
  \item[Mutacja] losowa modyfikacja osobnika.
  \item[Elityzm] przenoszenie najlepszych osobników bez zmian do kolejnego pokolenia.
\end{description}

\subsubsection{Ogólny schemat}

\begin{enumerate}
  \item Zainicjalizuj populację, losowo albo heurystycznie.
  \item Oceń każdego osobnika funkcją dopasowania.
  \item Dopóki nie spełniono warunku stopu:
  \begin{enumerate}
    \item wybierz rodziców,
    \item zastosuj krzyżowanie,
    \item zastosuj mutację,
    \item oceń potomków,
    \item utwórz nowe pokolenie,
    \item zachowaj elity, jeśli stosujesz elityzm.
  \end{enumerate}
  \item Zwróć najlepszego znalezionego osobnika.
\end{enumerate}

Warunek stopu może zależeć od liczby pokoleń, budżetu czasu, liczby ewaluacji funkcji celu, braku poprawy albo osiągnięcia wymaganej jakości.

\subsubsection{Reprezentacja}

Reprezentacja jest kluczowa. Powinna:
\begin{itemize}
  \item kodować wszystkie istotne rozwiązania,
  \item ograniczać liczbę rozwiązań niedopuszczalnych,
  \item pozwalać operatorom tworzyć sensownych potomków,
  \item zachowywać strukturę problemu.
\end{itemize}

Przykłady:
\begin{itemize}
  \item bitstring dla wyboru podzbioru cech,
  \item wektor rzeczywisty dla optymalizacji ciągłej,
  \item permutacja dla TSP i harmonogramowania,
  \item drzewo dla programowania genetycznego.
\end{itemize}

Jeśli reprezentacja generuje rozwiązania niedopuszczalne, trzeba stosować kary, naprawę, dekoder albo specjalne operatory zachowujące dopuszczalność.

\subsubsection{Selekcja}

Selekcja zwiększa prawdopodobieństwo reprodukcji lepszych osobników, ale nie powinna zbyt szybko usuwać różnorodności.

\begin{description}
  \item[Ruletka] prawdopodobieństwo wyboru proporcjonalne do dopasowania. Wrażliwa na skalowanie dopasowania.
  \item[Selekcja rankingowa] osobniki sortuje się według jakości, a prawdopodobieństwa zależą od rangi. Jest stabilniejsza.
  \item[Turniej] losuje się kilka osobników i wybiera najlepszego. Rozmiar turnieju kontroluje presję selekcyjną.
  \item[Elityzm] najlepsze osobniki przechodzą do następnego pokolenia, co zapobiega utracie najlepszego znalezionego rozwiązania.
\end{description}

\subsubsection{Krzyżowanie i mutacja}

Dla bitstringów używa się krzyżowania jednopunktowego, wielopunktowego albo jednorodnego. Dla wektorów rzeczywistych można stosować średnie arytmetyczne, BLX-$\alpha$, SBX albo rekombinację liniową. Dla permutacji potrzebne są operatory zachowujące permutację, np. PMX, OX, CX.

Mutacja zapewnia eksplorację. Przykłady:
\begin{itemize}
  \item odwrócenie bitu,
  \item dodanie szumu Gaussowskiego,
  \item zamiana dwóch pozycji w permutacji,
  \item odwrócenie fragmentu trasy,
  \item mutacja poddrzewa w programowaniu genetycznym.
\end{itemize}

\subsubsection{Eksploracja i eksploatacja}

Algorytm ewolucyjny musi równoważyć:
\begin{description}
  \item[Eksplorację] przeszukiwanie nowych regionów przestrzeni.
  \item[Eksploatację] intensywne ulepszanie obiecujących regionów.
\end{description}

Zbyt duża presja selekcyjna albo zbyt mała mutacja prowadzą do przedwczesnej zbieżności. Zbyt duża losowość powoduje zachowanie podobne do losowego przeszukiwania. Różnorodność można utrzymywać przez mutację, niching, crowding, fitness sharing, restart albo populacje wyspowe.

\subsubsection{Parametry}

Najważniejsze parametry:
\begin{itemize}
  \item rozmiar populacji,
  \item prawdopodobieństwo krzyżowania,
  \item prawdopodobieństwo mutacji,
  \item siła mutacji,
  \item presja selekcyjna,
  \item liczba elit,
  \item kryterium stopu.
\end{itemize}

Parametry silnie wpływają na wynik, dlatego należy je stroić na walidacyjnych instancjach i raportować w eksperymentach.

\subsubsection{Co powiedzieć na obronie}

Algorytm ewolucyjny utrzymuje populację rozwiązań, ocenia je funkcją dopasowania, wybiera lepsze osobniki do reprodukcji i tworzy nowe rozwiązania przez krzyżowanie oraz mutację. Selekcja daje presję w stronę jakości, mutacja i różnorodność zapewniają eksplorację, a elityzm chroni najlepsze rozwiązania. Kluczowe są reprezentacja i dobór operatorów zgodnych z problemem.

\subsection{Porównanie algorytmu ewolucyjnego z metodami optymalizacji w $\R^n$}

Algorytmy ewolucyjne i klasyczne metody optymalizacji ciągłej rozwiązują podobny problem poszukiwania ekstremum, ale mają inne założenia, koszty i gwarancje.

\subsubsection{Algorytm ewolucyjny}

Algorytm ewolucyjny:
\begin{itemize}
  \item jest populacyjny,
  \item zwykle nie wymaga gradientu,
  \item dobrze radzi sobie z funkcjami niegładkimi, zaszumionymi i wielomodalnymi,
  \item łatwo obsługuje reprezentacje dyskretne i mieszane,
  \item może znajdować dobre rozwiązania globalnie, ale bez gwarancji optimum,
  \item wymaga wielu ewaluacji funkcji celu.
\end{itemize}

Jakość wyniku jest losowa i zależy od parametrów oraz budżetu obliczeniowego. Wyniki powinno się raportować statystycznie.

\subsubsection{Metoda największego spadku}

Gradient descent:
\[
x_{k+1}=x_k-\alpha_k\nabla f(x_k)
\]
jest metodą lokalną, deterministyczną przy ustalonym punkcie startowym i kroku. Dla funkcji gładkich wypukłych ma gwarancje zbieżności, a dla silnie wypukłych zbieżność liniową. Wymaga gradientu i jest wrażliwa na skalowanie zmiennych oraz uwarunkowanie. W funkcjach niewypukłych może utknąć w minimum lokalnym albo punkcie siodłowym.

\subsubsection{Metoda Neldera-Meada}

Nelder-Mead jest bezpochodną metodą lokalną. Utrzymuje sympleks $n+1$ punktów i wykonuje operacje odbicia, ekspansji, kontrakcji i redukcji. Dobrze działa w małych wymiarach i dla funkcji bez gradientu, ale w dużych wymiarach traci efektywność. Nie jest populacyjny w sensie ewolucyjnym, choć utrzymuje wiele punktów. Ma słabsze gwarancje zbieżności niż metody gradientowe w problemach gładkich wypukłych.

\subsubsection{Metody quasi-Newtonowskie}

BFGS i L-BFGS korzystają z gradientu i przybliżają krzywiznę funkcji. Dla gładkich problemów są zwykle znacznie szybsze niż gradient descent. W problemach wypukłych mogą dawać bardzo wysoką dokładność. Wymagają jednak gradientu lub jego dobrego przybliżenia, a dla funkcji silnie zaszumionych albo niedyskretnej reprezentacji mogą być nieodpowiednie.

\subsubsection{Porównanie jakości i sposobu działania}

\begin{longtable}{p{0.22\textwidth}p{0.35\textwidth}p{0.35\textwidth}}
\toprule
Kryterium & Algorytm ewolucyjny & Metody wypukłe/lokalne \\
\midrule
Informacja o funkcji & wartości funkcji, czasem ograniczenia & gradient, czasem Hessian albo pochodne \\
Charakter & globalna metaheurystyka populacyjna & lokalna metoda analityczna \\
Wypukłość & nie wymaga wypukłości & wypukłość daje silne gwarancje \\
Gładkość & nie wymaga gładkości & zwykle wymaga gładkości \\
Koszt & wiele ewaluacji funkcji celu & mniej iteracji, ale każda może wymagać gradientu/Hessianu \\
Wynik & losowy, zależny od budżetu & deterministyczny przy ustalonych parametrach \\
Gwarancje & zwykle słabe formalnie & silne dla problemów wypukłych \\
Duży wymiar & często trudny bez specjalizacji & L-BFGS/SGD skalują się dobrze \\
\bottomrule
\end{longtable}

\subsubsection{Kiedy wybrać którą metodę}

Algorytm ewolucyjny warto wybrać, gdy:
\begin{itemize}
  \item funkcja jest czarnoskrzynkowa,
  \item nie ma gradientu,
  \item przestrzeń jest dyskretna albo mieszana,
  \item problem jest wielomodalny,
  \item ewaluacje można równoleglić,
  \item wystarczy dobre rozwiązanie, niekoniecznie dowód optymalności.
\end{itemize}

Metody gradientowe/quasi-Newtonowskie warto wybrać, gdy:
\begin{itemize}
  \item funkcja jest gładka,
  \item gradient jest dostępny,
  \item problem jest wypukły albo lokalna dokładność jest ważna,
  \item wymiar jest duży,
  \item liczba ewaluacji funkcji celu ma być mała.
\end{itemize}

Często stosuje się hybrydy: algorytm ewolucyjny znajduje obiecujący region, a metoda lokalna dopracowuje rozwiązanie. Takie podejście nazywa się czasem algorytmem memetycznym.

\subsubsection{Co powiedzieć na obronie}

Algorytm ewolucyjny jest globalną, losową, populacyjną heurystyką bez wymogu gradientu, dobrą dla funkcji trudnych, niegładkich i wielomodalnych. Metody gradientowe, Nelder-Mead i quasi-Newton są bardziej lokalne; dla problemów wypukłych i gładkich dają znacznie lepsze gwarancje oraz dokładność. Ewolucja jest bardziej uniwersalna, ale kosztowna i mniej przewidywalna; metody wypukłe są szybsze i lepiej uzasadnione, gdy spełnione są ich założenia.

\subsection{Metody symulacji Monte Carlo w grach}

Metody Monte Carlo w grach używają losowych symulacji rozgrywki do oceny ruchów. Są szczególnie użyteczne, gdy przestrzeń stanów jest ogromna, trudno zbudować dobrą funkcję oceny, a pełne przeszukiwanie minimax jest niemożliwe.

\subsubsection{Podstawowa idea}

Dla danego stanu gry i możliwego ruchu wykonujemy wiele losowych dogrywek do końca gry albo do ustalonej głębokości. Wynik ruchu oceniamy średnią nagrodą:
\[
\hat V(a)=\frac{1}{N}\sum_{i=1}^{N} z_i,
\]
gdzie $z_i$ jest wynikiem $i$-tej symulacji po ruchu $a$. Ruch z największą średnią jest wybierany.

Zaletą jest prostota i brak konieczności ręcznego projektowania skomplikowanej heurystyki. Wadą jest wariancja i koszt wielu symulacji.

\subsubsection{Monte Carlo Tree Search}

MCTS łączy symulacje Monte Carlo z budową drzewa przeszukiwania. Najczęściej opisuje się cztery fazy:
\begin{enumerate}
  \item selekcja - przechodzimy od korzenia do liścia, wybierając dzieci według reguły równoważącej eksplorację i eksploatację,
  \item ekspansja - dodajemy nowy węzeł,
  \item symulacja - wykonujemy dogrywkę, np. losową albo heurystyczną,
  \item propagacja wsteczna - aktualizujemy statystyki odwiedzonych węzłów.
\end{enumerate}

W każdym węźle przechowuje się zwykle liczbę wizyt $N(s)$ i sumę albo średnią nagród $Q(s,a)$.

\subsubsection{UCT}

Popularną regułą selekcji jest UCT, oparta na upper confidence bound:
\[
a=\argmax_a\left(\bar X(s,a)+C\sqrt{\frac{\ln N(s)}{N(s,a)}}\right).
\]
Pierwszy składnik promuje ruchy o wysokiej średniej wygranej, drugi promuje ruchy rzadko badane. Parametr $C$ kontroluje eksplorację.

\subsubsection{Zastosowania}

MCTS był przełomowy w grach takich jak Go, gdzie klasyczne minimax z ręczną funkcją oceny długo było niewystarczające. Stosuje się go też w grach planszowych, grach z niepełną informacją w wariantach determinization albo information set MCTS, planowaniu robotów, grach wideo i problemach decyzyjnych.

\subsubsection{Ulepszenia}

\begin{itemize}
  \item rollout policy zamiast całkowicie losowych dogrywek,
  \item progressive widening dla dużej liczby akcji,
  \item transposition tables,
  \item RAVE/AMAF,
  \item priory z sieci neuronowej,
  \item value network zastępująca część symulacji,
  \item parallel MCTS.
\end{itemize}

AlphaGo i AlphaZero łączyły MCTS z sieciami neuronowymi: sieć polityki sugerowała obiecujące ruchy, a sieć wartości oceniała pozycje.

\subsubsection{Co powiedzieć na obronie}

Monte Carlo w grach ocenia ruchy przez losowe symulacje. MCTS buduje asymetryczne drzewo tam, gdzie symulacje wskazują obiecujące warianty. Kluczowe fazy to selekcja, ekspansja, symulacja i backpropagation. UCT równoważy eksploatację ruchów dobrych i eksplorację słabo poznanych. Metody te są skuteczne przy ogromnych drzewach gry i słabych heurystykach.

\subsection{Pojęcie agenta w SI: opis formalny i własności}

Agent w sztucznej inteligencji to system, który postrzega środowisko przez sensory i oddziałuje na nie przez aktuatory. Formalnie agent realizuje odwzorowanie historii percepcji na akcje:
\[
f:P^*\to A,
\]
gdzie $P^*$ jest zbiorem możliwych historii perceptów, a $A$ zbiorem akcji. Często zamiast historii używa się stanu wewnętrznego, który jest jej skompresowaną reprezentacją.

\subsubsection{Agent i środowisko}

Cykl działania:
\[
\text{percepcja}\to\text{aktualizacja stanu}\to\text{wybór akcji}\to\text{zmiana środowiska}.
\]
Środowisko może być:
\begin{itemize}
  \item w pełni obserwowalne albo częściowo obserwowalne,
  \item deterministyczne albo stochastyczne,
  \item epizodyczne albo sekwencyjne,
  \item statyczne albo dynamiczne,
  \item dyskretne albo ciągłe,
  \item jednoagentowe albo wieloagentowe,
  \item znane albo nieznane.
\end{itemize}

\subsubsection{Racjonalność}

Agent racjonalny wybiera akcję maksymalizującą oczekiwaną miarę skuteczności na podstawie dostępnej wiedzy i perceptów. Racjonalność nie oznacza wszechwiedzy. Agent może podjąć decyzję błędną ex post, jeśli była najlepsza ex ante przy dostępnej informacji.

Formalnie można mówić o maksymalizacji oczekiwanej użyteczności:
\[
a^*=\argmax_a \mathbb{E}[U(\text{wynik})\mid \text{historia},a].
\]

\subsubsection{Typy agentów}

\begin{description}
  \item[Prosty agent reaktywny] wybiera akcję na podstawie bieżącej percepcji, reguł typu warunek-akcja.
  \item[Agent ze stanem] utrzymuje stan wewnętrzny opisujący nieobserwowalne aspekty świata.
  \item[Agent celowy] wybiera akcje prowadzące do celu.
  \item[Agent użytecznościowy] porównuje różne wyniki przez funkcję użyteczności.
  \item[Agent uczący się] poprawia działanie na podstawie doświadczenia.
\end{description}

\subsubsection{Własności agentów}

W systemach wieloagentowych często wymienia się:
\begin{itemize}
  \item autonomię - agent działa bez bezpośredniego sterowania,
  \item reaktywność - reaguje na zmiany środowiska,
  \item proaktywność - inicjuje działania dla osiągania celów,
  \item społeczność - komunikuje się i współpracuje z innymi agentami,
  \item adaptacyjność - zmienia zachowanie na podstawie doświadczenia,
  \item racjonalność - dobiera akcje zgodnie z miarą skuteczności,
  \item ograniczoną racjonalność - działa przy ograniczonych zasobach obliczeniowych.
\end{itemize}

\subsubsection{Agent w uczeniu ze wzmocnieniem}

W RL agent w stanie $S_t$ wybiera akcję $A_t$, środowisko zwraca nagrodę $R_{t+1}$ i kolejny stan $S_{t+1}$. Celem jest maksymalizacja oczekiwanego zdyskontowanego zwrotu:
\[
G_t=R_{t+1}+\gamma R_{t+2}+\gamma^2R_{t+3}+\dots.
\]
Polityka $\pi(a|s)$ opisuje zachowanie agenta.

\subsubsection{Co powiedzieć na obronie}

Agent to system postrzegający środowisko i wykonujący akcje. Formalnie można go opisać jako funkcję z historii perceptów do akcji. Agent racjonalny maksymalizuje oczekiwaną miarę skuteczności przy dostępnej wiedzy. Ważne własności to autonomia, reaktywność, proaktywność, zdolność komunikacji i uczenie się.

\subsection{Inteligencja rojowa: idea i przykłady realizacji}

Inteligencja rojowa opisuje rozwiązywanie problemów przez wiele prostych jednostek, które lokalnie oddziałują ze sobą i środowiskiem, a globalnie tworzą złożone, adaptacyjne zachowanie. Inspiracją są mrówki, pszczoły, ptaki, ryby i inne systemy kolektywne.

\subsubsection{Cechy systemów rojowych}

\begin{itemize}
  \item brak centralnego sterowania,
  \item lokalne reguły zachowania,
  \item komunikacja bezpośrednia albo przez środowisko,
  \item samoorganizacja,
  \item redundancja i odporność,
  \item równoległość,
  \item emergencja, czyli powstawanie globalnego wzorca z lokalnych interakcji.
\end{itemize}

Stigmergia oznacza komunikację przez ślady w środowisku. Klasyczny przykład to feromony mrówek.

\subsubsection{Ant Colony Optimization}

ACO rozwiązuje problemy kombinatoryczne przez sztuczne mrówki budujące rozwiązania krok po kroku. Prawdopodobieństwo wyboru elementu zależy od feromonu i heurystyki lokalnej:
\[
p_{ij}\propto \tau_{ij}^{\alpha}\eta_{ij}^{\beta},
\]
gdzie $\tau_{ij}$ to poziom feromonu, a $\eta_{ij}$ to atrakcyjność heurystyczna, np. odwrotność odległości w TSP.

Po zbudowaniu rozwiązań feromon jest aktualizowany:
\[
\tau_{ij}\leftarrow (1-\rho)\tau_{ij}+\Delta\tau_{ij}.
\]
Parowanie feromonu zapobiega zbyt szybkiemu zablokowaniu na jednym rozwiązaniu.

\subsubsection{Particle Swarm Optimization}

PSO działa w przestrzeni ciągłej. Cząstka ma pozycję $x_i$ i prędkość $v_i$. Aktualizacja:
\[
v_i\leftarrow \omega v_i+c_1r_1(p_i-x_i)+c_2r_2(g-x_i),
\]
\[
x_i\leftarrow x_i+v_i.
\]
Składniki oznaczają: bezwładność, przyciąganie do najlepszego własnego punktu $p_i$ i przyciąganie do najlepszego punktu roju $g$. PSO jest proste i skuteczne dla wielu problemów ciągłych, ale może przedwcześnie zbiegać.

\subsubsection{Artificial Bee Colony}

ABC inspiruje się zachowaniem pszczół. Rozwiązania są źródłami pożywienia. Występują pszczoły zatrudnione, obserwatorki i zwiadowczynie. Dobre źródła są eksploatowane, słabe porzucane, a zwiadowczynie szukają nowych regionów.

\subsubsection{Boids i modele stadne}

W symulacji stad ptaków lub ryb globalny ruch powstaje z trzech lokalnych reguł:
\begin{itemize}
  \item separacja - unikaj kolizji,
  \item alignment - dopasuj kierunek do sąsiadów,
  \item cohesion - trzymaj się grupy.
\end{itemize}
To przykład emergencji: z prostych reguł lokalnych powstaje złożone zachowanie zbiorowe.

\subsubsection{Zastosowania}

\begin{itemize}
  \item TSP i optymalizacja tras,
  \item harmonogramowanie,
  \item dobór parametrów modeli,
  \item robotyka rojowa,
  \item routing w sieciach,
  \item optymalizacja funkcji ciągłych,
  \item klastrowanie i eksploracja danych.
\end{itemize}

\subsubsection{Co powiedzieć na obronie}

Inteligencja rojowa polega na tym, że wiele prostych jednostek działających lokalnie tworzy globalnie inteligentne zachowanie. Nie ma centralnego sterowania, ważna jest samoorganizacja i komunikacja lokalna. Przykłady to ACO dla problemów kombinatorycznych, PSO dla optymalizacji ciągłej, ABC oraz modele stadne.

\subsection{Programowanie w logice na przykładzie Prologu}

Programowanie logiczne opisuje problem przez fakty i reguły, a obliczenie polega na logicznym wnioskowaniu. Prolog jest klasycznym językiem programowania logicznego opartym na klauzulach Horna, unifikacji i rezolucji SLD.

\subsubsection{Fakty, reguły i zapytania}

Fakt:
\begin{lstlisting}[language=Prolog]
parent(jan, anna).
parent(anna, piotr).
\end{lstlisting}

Reguła:
\begin{lstlisting}[language=Prolog]
grandparent(X, Z) :- parent(X, Y), parent(Y, Z).
\end{lstlisting}

Zapytanie:
\begin{lstlisting}[language=Prolog]
?- grandparent(jan, piotr).
\end{lstlisting}

Regułę czytamy: $X$ jest dziadkiem/babcią $Z$, jeśli istnieje $Y$, takie że $X$ jest rodzicem $Y$ i $Y$ jest rodzicem $Z$.

\subsubsection{Semantyka klauzul Horna}

Klauzula Horna ma co najwyżej jeden literał pozytywny. Reguła Prologa:
\[
H\;:-\;B_1,\dots,B_k
\]
odpowiada formule:
\[
B_1\wedge\dots\wedge B_k\to H.
\]
Fakt to reguła z pustym ciałem. Zapytanie jest celem, którego Prolog próbuje dowieść.

\subsubsection{Unifikacja}

Unifikacja dopasowuje termy przez podstawienia. Dla zapytania:
\begin{lstlisting}[language=Prolog]
?- parent(jan, X).
\end{lstlisting}
i faktu:
\begin{lstlisting}[language=Prolog]
parent(jan, anna).
\end{lstlisting}
unifikator to $\{X/anna\}$. Prolog zwróci $X=anna$.

\subsubsection{Rezolucja SLD i backtracking}

Prolog bierze pierwszy cel z listy celów, szuka pierwszej pasującej klauzuli, unifikuje, zastępuje cel ciałem reguły i kontynuuje. Jeśli później dojdzie do porażki, cofa się do ostatniego punktu wyboru i próbuje kolejnej klauzuli. To backtracking.

Kolejność reguł i kolejność celów w ciele ma znaczenie operacyjne, mimo że deklaratywnie mogą być równoważne. Źle ustawiona rekursja może powodować zapętlenie.

\subsubsection{Rekursja}

Przykład przodka:
\begin{lstlisting}[language=Prolog]
ancestor(X, Z) :- parent(X, Z).
ancestor(X, Z) :- parent(X, Y), ancestor(Y, Z).
\end{lstlisting}
Pierwsza reguła jest przypadkiem bazowym, druga rekurencyjnym.

\subsubsection{Negacja jako porażka}

Prolog często używa negation as failure:
\begin{lstlisting}[language=Prolog]
not_p(X) :- \+ p(X).
\end{lstlisting}
Znaczy to: uznaj $\neg p(X)$, jeśli nie udało się dowieść $p(X)$. To nie jest klasyczna negacja logiczna; opiera się na założeniu zamkniętego świata. Jeśli czegoś nie ma w bazie i nie da się tego wywnioskować, traktujemy to jako fałsz.

\subsubsection{Cut}

Operator cut `!' ogranicza backtracking. Może poprawiać wydajność albo wymuszać deterministyczny wybór, ale osłabia deklaratywną czytelność programu. Cut zielony nie zmienia znaczenia logicznego, tylko przyspiesza. Cut czerwony zmienia odpowiedzi programu i wymaga ostrożności.

\subsubsection{Co powiedzieć na obronie}

Prolog reprezentuje wiedzę przez fakty i reguły będące klauzulami Horna. Obliczenie polega na dowodzeniu zapytań przez unifikację, rezolucję SLD i backtracking. Program ma charakter deklaratywny, ale kolejność reguł i celów wpływa na działanie. Ważne pojęcia to unifikacja, rekursja, negacja jako porażka i cut.

\subsection{Generowanie reguł minimalnych z bazy wiedzy w systemach eksperckich}

System ekspertowy używa bazy wiedzy i mechanizmu wnioskowania do rozwiązywania problemów z danej dziedziny. Wiedza często ma postać reguł:
\[
IF\; warunki \; THEN\; decyzja.
\]
Reguły minimalne są regułami, z których nie można usunąć żadnego warunku bez utraty zdolności rozróżniania decyzji albo bez zmiany pokrycia zgodnego z założonym kryterium.

\subsubsection{System informacyjny i decyzyjny}

Dane można przedstawić jako tabelę decyzyjną:
\[
S=(U,A\cup\{d\}),
\]
gdzie $U$ jest zbiorem obiektów, $A$ zbiorem atrybutów warunkowych, a $d$ atrybutem decyzyjnym. Reguła ma postać:
\[
(a_1=v_1)\wedge\dots\wedge(a_k=v_k)\to(d=c).
\]

\subsubsection{Minimalność reguły}

Reguła jest minimalna, jeśli:
\begin{itemize}
  \item jest poprawna albo wystarczająco wiarygodna według przyjętej miary,
  \item konkluzja pozostaje określona przez warunki,
  \item żaden warunek z poprzednika nie jest redundantny.
\end{itemize}

Minimalność nie zawsze oznacza najkrótszą możliwą regułę globalnie. Może oznaczać minimalność względem inkluzji zbioru warunków.

\subsubsection{Podejście przez zbiory przybliżone}

W teorii zbiorów przybliżonych obiekty są nierozróżnialne względem zbioru atrybutów $B$, jeśli mają te same wartości wszystkich atrybutów z $B$. Redukt to minimalny podzbiór atrybutów zachowujący zdolność klasyfikacyjną całego zbioru atrybutów. Core to część wspólna wszystkich reduktów, czyli atrybuty konieczne.

Generowanie reguł:
\begin{enumerate}
  \item zbudować klasy decyzyjne,
  \item wyznaczyć relacje nierozróżnialności,
  \item znaleźć redukty lokalne albo globalne,
  \item utworzyć reguły z minimalnych kombinacji atrybut-wartość,
  \item usunąć warunki redundantne,
  \item ocenić wsparcie, dokładność i pokrycie.
\end{enumerate}

\subsubsection{Macierz rozróżnialności}

Dla par obiektów o różnych decyzjach buduje się zbiory atrybutów, na których obiekty się różnią. Aby reguła odróżniała obiekt od obiektów z inną decyzją, musi zawierać co najmniej jeden atrybut z każdego takiego zbioru. Problem sprowadza się do znalezienia minimalnych hitting sets, co może być kosztowne obliczeniowo.

\subsubsection{Miary reguł}

Dla reguły $r: P\to C$:
\begin{description}
  \item[Wsparcie] liczba obiektów spełniających $P$ i $C$.
  \item[Pokrycie] jaka część klasy decyzyjnej jest pokryta przez regułę.
  \item[Dokładność] odsetek obiektów spełniających $P$, które mają decyzję $C$:
  \[
  acc(r)=\frac{|P\cap C|}{|P|}.
  \]
  \item[Confidence] analog dokładności w regułach asocjacyjnych.
  \item[Lift] porównanie z częstością klasy bazowej.
\end{description}

W danych zaszumionych wymóg stuprocentowej dokładności może prowadzić do zbyt specyficznych reguł. Dopuszcza się reguły przybliżone z minimalnym wsparciem i minimalną dokładnością.

\subsubsection{Usuwanie redundancji}

Regułę minimalizuje się przez testowanie każdego warunku:
\begin{enumerate}
  \item usuń warunek,
  \item sprawdź, czy reguła nadal spełnia kryterium jakości,
  \item jeśli tak, warunek był redundantny,
  \item powtarzaj do punktu stałego.
\end{enumerate}

Usuwanie pojedynczych warunków nie zawsze daje globalnie najkrótszą regułę, ale jest praktyczne. Dokładniejsze metody przeszukują przestrzeń podzbiorów warunków.

\subsubsection{Związek z drzewami decyzyjnymi}

Reguły można generować z drzewa decyzyjnego: każda ścieżka od korzenia do liścia tworzy regułę. Następnie reguły można przycinać, usuwając warunki, które nie pogarszają jakości na danych walidacyjnych. Takie reguły są czytelne, choć zależą od heurystyki budowy drzewa.

\subsubsection{Co powiedzieć na obronie}

Reguły minimalne to reguły decyzyjne bez zbędnych warunków. Generuje się je z tabeli decyzyjnej przez analizę rozróżnialności obiektów, redukty atrybutów, zbiory przybliżone albo drzewa decyzyjne. Minimalizacja polega na usuwaniu warunków, które nie są potrzebne do zachowania decyzji lub jakości reguły. Reguły ocenia się przez wsparcie, dokładność i pokrycie.

\section{Podstawy przetwarzania danych}

\subsection{Metody redukcji wymiaru danych}

Redukcja wymiaru polega na zastąpieniu danych opisanych wieloma zmiennymi przez reprezentację o mniejszej liczbie wymiarów przy możliwie małej utracie istotnej informacji. Stosuje się ją dla wizualizacji, przyspieszenia algorytmów, usuwania szumu, walki z przekleństwem wymiarowości i poprawy generalizacji.

\subsubsection{Redukcja wymiaru a selekcja cech}

Redukcja wymiaru tworzy nowe zmienne będące transformacjami cech oryginalnych. Selekcja cech wybiera podzbiór istniejących cech. PCA tworzy nowe składowe liniowe; selekcja może wybrać np. tylko wiek, dochód i liczbę transakcji z oryginalnej tabeli.

\subsubsection{PCA}

Principal Component Analysis znajduje ortogonalne kierunki maksymalnej wariancji danych. Po centrowaniu macierzy danych $X$ szukamy kierunku $w$, który maksymalizuje:
\[
\operatorname{Var}(Xw)
\]
przy ograniczeniu $\|w\|=1$. Rozwiązaniem są wektory własne macierzy kowariancji:
\[
\Sigma=\frac{1}{n-1}X^TX.
\]
Pierwsza składowa główna wyjaśnia największą wariancję, druga największą wariancję w kierunku ortogonalnym do pierwszej itd.

Zalety PCA:
\begin{itemize}
  \item prostota,
  \item szybkość,
  \item dekorrelacja zmiennych,
  \item możliwość interpretacji przez wariancję wyjaśnioną.
\end{itemize}
Wady:
\begin{itemize}
  \item liniowość,
  \item wrażliwość na skalowanie cech,
  \item skupienie na wariancji, niekoniecznie na informacji predykcyjnej,
  \item ograniczona interpretowalność składowych.
\end{itemize}

\subsubsection{SVD}

Singular Value Decomposition rozkłada macierz:
\[
X=U\Sigma V^T.
\]
Przycięcie do $k$ największych wartości osobliwych daje najlepsze przybliżenie rzędu $k$ w normie Frobeniusa. SVD jest podstawą PCA i metod latent semantic analysis.

\subsubsection{LDA jako redukcja nadzorowana}

Linear Discriminant Analysis szuka kierunków dobrze rozdzielających klasy. Maksymalizuje stosunek wariancji międzyklasowej do wewnątrzklasowej. Dla $C$ klas redukuje do co najwyżej $C-1$ wymiarów. W przeciwieństwie do PCA używa etykiet.

\subsubsection{Metody nieliniowe}

\begin{description}
  \item[Kernel PCA] wykonuje PCA w przestrzeni cech zadanej przez kernel.
  \item[Isomap] zachowuje odległości geodezyjne na rozmaitości.
  \item[LLE] zakłada, że punkt można odtworzyć z lokalnych sąsiadów.
  \item[t-SNE] służy głównie do wizualizacji, zachowuje lokalne podobieństwa, ale zniekształca globalne odległości.
  \item[UMAP] szybka metoda wizualizacji i redukcji, często lepiej zachowuje strukturę globalną niż t-SNE, ale nadal wymaga ostrożnej interpretacji.
  \item[Autoenkodery] sieci neuronowe uczące kodu niskowymiarowego przez rekonstrukcję wejścia.
\end{description}

\subsubsection{Dobór liczby wymiarów}

Można użyć:
\begin{itemize}
  \item progu wariancji wyjaśnionej, np. 95\% w PCA,
  \item wykresu osypiska,
  \item walidacji krzyżowej dla jakości modelu predykcyjnego,
  \item metryk rekonstrukcji,
  \item stabilności reprezentacji.
\end{itemize}

\subsubsection{Co powiedzieć na obronie}

Redukcja wymiaru tworzy niskowymiarową reprezentację danych. PCA jest podstawową metodą liniową maksymalizującą wariancję, SVD daje jej stabilną implementację, LDA używa etykiet, a metody nieliniowe takie jak t-SNE, UMAP i autoenkodery wykrywają bardziej złożone struktury. Trzeba pamiętać o skalowaniu, utracie informacji i walidacji liczby wymiarów.

\subsection{Metody selekcji cech}

Selekcja cech wybiera podzbiór oryginalnych zmiennych. Celem jest poprawa generalizacji, redukcja kosztu pomiaru, przyspieszenie modelu, interpretowalność i zmniejszenie ryzyka overfittingu.

\subsubsection{Filtry}

Metody filtrujące oceniają cechy niezależnie od konkretnego modelu albo z użyciem prostych statystyk:
\begin{itemize}
  \item wariancja cechy,
  \item korelacja z etykietą,
  \item test $\chi^2$ dla cech kategorycznych,
  \item ANOVA F-test,
  \item mutual information,
  \item współczynnik korelacji rang Spearmana,
  \item usuwanie jednej z silnie skorelowanych cech.
\end{itemize}

Filtry są szybkie i skalowalne, ale mogą ignorować interakcje między cechami.

\subsubsection{Wrappery}

Wrapper używa modelu jako czarnej skrzynki i ocenia podzbiory cech przez walidację:
\begin{itemize}
  \item forward selection - zaczynamy od pustego zbioru i dodajemy cechy,
  \item backward elimination - zaczynamy od wszystkich i usuwamy,
  \item recursive feature elimination,
  \item przeszukiwanie heurystyczne, np. algorytm genetyczny.
\end{itemize}

Wrappery mogą dawać lepsze podzbiory dla konkretnego modelu, ale są kosztowne i podatne na przeuczenie, jeśli walidacja jest źle zorganizowana.

\subsubsection{Metody wbudowane}

Embedded methods wybierają cechy w trakcie uczenia modelu:
\begin{itemize}
  \item Lasso, czyli regularyzacja $L_1$, zeruje część współczynników,
  \item Elastic Net łączy $L_1$ i $L_2$,
  \item drzewa i lasy losowe dają importance cech,
  \item gradient boosting ma miary gain/split importance,
  \item modele liniowe z karami grupowymi wybierają grupy cech.
\end{itemize}

Trzeba uważać na interpretację importance w drzewach: cechy o wielu możliwych podziałach albo silnie skorelowane mogą być faworyzowane lub dzielić znaczenie.

\subsubsection{Stability selection}

Stability selection wielokrotnie losuje podpróbki danych, wykonuje selekcję i wybiera cechy pojawiające się stabilnie. Pomaga odróżnić cechy rzeczywiście istotne od przypadkowych, zwłaszcza gdy liczba cech jest duża.

\subsubsection{Selekcja a przeciek danych}

Selekcja cech musi być wykonywana wewnątrz pętli walidacji. Błąd: wybrać cechy na całym zbiorze, a potem zrobić cross-validation. Informacja ze zbioru testowego przecieka wtedy do modelu i wynik jest zawyżony.

\subsubsection{Co powiedzieć na obronie}

Selekcja cech wybiera podzbiór oryginalnych zmiennych. Filtry są szybkie i niezależne od modelu, wrappery oceniają podzbiory przez model i walidację, a metody wbudowane wybierają cechy podczas uczenia, np. Lasso albo drzewa. Najważniejsze jest unikanie przecieku danych i ocena stabilności wyboru.

\subsection{Metody postępowania z brakami w danych}

Braki danych mogą zniekształcić analizę i modelowanie. Wybór metody zależy od mechanizmu braków, udziału braków, typu zmiennej i modelu docelowego.

\subsubsection{Mechanizmy braków}

\begin{description}
  \item[MCAR] Missing Completely At Random - brak niezależny od danych obserwowanych i nieobserwowanych. Usunięcie wierszy nie wprowadza biasu, ale zmniejsza moc statystyczną.
  \item[MAR] Missing At Random - brak zależy od danych obserwowanych, ale nie od brakującej wartości po uwzględnieniu obserwowanych zmiennych. Można go modelować.
  \item[MNAR] Missing Not At Random - brak zależy od samej brakującej wartości albo ukrytej zmiennej. Najtrudniejszy przypadek.
\end{description}

Przykład MNAR: osoby z bardzo wysokim dochodem częściej nie podają dochodu.

\subsubsection{Usuwanie danych}

\begin{itemize}
  \item listwise deletion - usunięcie wierszy z jakimkolwiek brakiem,
  \item pairwise deletion - użycie dostępnych par obserwacji dla danej statystyki,
  \item usuwanie cech z bardzo dużym udziałem braków.
\end{itemize}

Usuwanie jest proste, ale może wprowadzać bias i marnować dane, jeśli braki nie są MCAR.

\subsubsection{Prosta imputacja}

\begin{itemize}
  \item średnia albo mediana dla zmiennych liczbowych,
  \item dominanta dla zmiennych kategorycznych,
  \item stała wartość specjalna, np. "unknown",
  \item imputacja grupowa, np. mediana w obrębie segmentu.
\end{itemize}

Mediana jest odporniejsza na odstające wartości niż średnia. Prosta imputacja zaniża wariancję i może osłabić zależności między zmiennymi.

\subsubsection{Imputacja modelowa}

\begin{itemize}
  \item k-nearest neighbors imputation,
  \item regresja predykcyjna,
  \item drzewa i lasy,
  \item EM dla modeli probabilistycznych,
  \item multiple imputation,
  \item autoenkodery albo modele generatywne.
\end{itemize}

Multiple imputation tworzy wiele uzupełnionych zbiorów, analizuje każdy i łączy wyniki, uwzględniając niepewność imputacji.

\subsubsection{Indykatory braków}

Czasem sam fakt braku jest informacyjny. Dodaje się cechę binarną:
\[
m_j =
\begin{cases}
1, & x_j\text{ było brakujące},\\
0, & \text{w przeciwnym razie}.
\end{cases}
\]
To bywa skuteczne przy MAR/MNAR, ale nie zastępuje analizy mechanizmu braków.

\subsubsection{Modele obsługujące braki}

Niektóre modele potrafią obsługiwać braki bez jawnej imputacji, np. wybrane implementacje drzew gradientowych mają kierunek domyślny dla braków. Nadal trzeba rozumieć, co model robi z brakami i czy zachowanie jest zgodne z problemem.

\subsubsection{Co powiedzieć na obronie}

Najpierw trzeba rozpoznać mechanizm braków: MCAR, MAR albo MNAR. Można usuwać wiersze lub cechy, imputować prostymi statystykami, imputować modelowo, stosować multiple imputation albo dodać indykatory braków. Imputację trzeba dopasować tylko na danych treningowych, aby uniknąć przecieku.

\subsection{Estymatory błędów}

Estymator błędu ocenia, jak dobrze model będzie działał na nowych danych. Błąd treningowy jest zwykle zbyt optymistyczny, bo model był dopasowany do tych danych.

\subsubsection{Podział train/validation/test}

Najprostszy schemat:
\begin{itemize}
  \item zbiór treningowy - uczenie parametrów,
  \item zbiór walidacyjny - dobór hiperparametrów i wybór modelu,
  \item zbiór testowy - końcowa bezstronna ocena.
\end{itemize}

Zbiór testowy powinien być użyty tylko raz na końcu. Wielokrotne podejmowanie decyzji na podstawie testu prowadzi do przeuczenia do testu.

\subsubsection{Walidacja krzyżowa}

W $k$-fold cross-validation dane dzieli się na $k$ części. Model trenuje się $k$ razy, za każdym razem na $k-1$ częściach, a testuje na pozostałej. Wyniki uśrednia się:
\[
\hat E_{CV}=\frac{1}{k}\sum_{i=1}^{k}E_i.
\]
Typowe wartości to $k=5$ albo $k=10$.

Leave-one-out CV to przypadek $k=n$. Ma mały bias, ale może mieć dużą wariancję i jest kosztowna.

\subsubsection{Stratyfikacja i dane zależne}

Dla klasyfikacji niezbalansowanej stosuje się stratified CV, aby proporcje klas były podobne w foldach. Dla szeregów czasowych nie wolno mieszać przyszłości z przeszłością; używa się walidacji kroczącej. Dla danych grupowych trzeba dzielić po grupach, aby obserwacje tej samej osoby, klienta albo urządzenia nie trafiały jednocześnie do treningu i testu.

\subsubsection{Bootstrap}

Bootstrap losuje próbki z powtórzeniami z danych treningowych. Obserwacje niewylosowane tworzą out-of-bag set. Bootstrap pozwala oceniać niepewność estymatora i stabilność modelu. W lasach losowych OOB error jest naturalnym estymatorem błędu.

\subsubsection{Macierz pomyłek i miary klasyfikacji}

Dla klasyfikacji binarnej:
\begin{itemize}
  \item TP - prawdziwie pozytywne,
  \item TN - prawdziwie negatywne,
  \item FP - fałszywie pozytywne,
  \item FN - fałszywie negatywne.
\end{itemize}

Miary:
\[
accuracy=\frac{TP+TN}{TP+TN+FP+FN},
\]
\[
precision=\frac{TP}{TP+FP},
\]
\[
recall=\frac{TP}{TP+FN},
\]
\[
F1=\frac{2\cdot precision\cdot recall}{precision+recall}.
\]
Accuracy jest mylące przy silnie niezbalansowanych klasach.

\subsubsection{Miary regresji}

\[
MAE=\frac{1}{n}\sum_i |y_i-\hat y_i|,
\]
\[
MSE=\frac{1}{n}\sum_i (y_i-\hat y_i)^2,
\]
\[
RMSE=\sqrt{MSE}.
\]
MAE jest odporniejsze na outliery, MSE/RMSE mocniej karzą duże błędy. $R^2$ mierzy część wariancji wyjaśnionej przez model.

\subsubsection{Bias i wariancja estymatora błędu}

Estymator błędu może być obciążony i mieć wariancję. Mały zbiór testowy daje dużą niepewność. Cross-validation zmniejsza zależność od pojedynczego podziału, ale wyniki foldów nie są całkowicie niezależne. Przy porównywaniu modeli warto stosować testy statystyczne albo przedziały ufności.

\subsubsection{Co powiedzieć na obronie}

Błąd modelu trzeba estymować na danych nieużytych do uczenia. Podstawowe metody to holdout, train/validation/test, walidacja krzyżowa, bootstrap i OOB. Miara błędu zależy od zadania: accuracy, precision, recall, F1, AUC dla klasyfikacji; MAE, MSE, RMSE dla regresji. Należy unikać przecieku danych i dobierać schemat walidacji do struktury danych.

\section{Uczenie ze wzmocnieniem}

\subsection{Uczenie ze wzmocnieniem: pasywne, aktywne i polityki}

Uczenie ze wzmocnieniem bada sytuację, w której agent uczy się przez interakcję ze środowiskiem. Agent wykonuje akcje, obserwuje następne stany i otrzymuje nagrody. Nie dostaje poprawnych etykiet akcji, lecz sygnał oceny skutków działania.

\subsubsection{Cykl interakcji}

W chwili $t$ agent jest w stanie $S_t$, wybiera akcję $A_t$, środowisko zwraca nagrodę $R_{t+1}$ i nowy stan $S_{t+1}$. Celem jest maksymalizacja oczekiwanego zwrotu:
\[
G_t=R_{t+1}+\gamma R_{t+2}+\gamma^2R_{t+3}+\dots
=\sum_{k=0}^{\infty}\gamma^k R_{t+k+1},
\]
gdzie $\gamma\in[0,1]$ jest współczynnikiem dyskontowym. Dla zadań epizodycznych suma kończy się w stanie terminalnym.

\subsubsection{Polityka}

Polityka opisuje sposób wyboru akcji:
\[
\pi(a|s)=P(A_t=a\mid S_t=s).
\]
Polityka deterministyczna przypisuje każdemu stanowi jedną akcję:
\[
\pi(s)=a.
\]
Polityka stochastyczna przypisuje rozkład prawdopodobieństwa na akcjach. Stochastyczność jest ważna dla eksploracji i dla środowisk częściowo obserwowalnych albo wieloagentowych.

\subsubsection{Pasywne uczenie ze wzmocnieniem}

W pasywnym RL polityka jest dana. Agent nie decyduje, jak ją ulepszyć, lecz ocenia jej jakość. Typowe zadanie: policy evaluation, czyli estymacja funkcji wartości $V^\pi(s)$ dla ustalonej polityki $\pi$.

Przykład: mamy strategię gry w blackjacka "dobieraj do 20, potem stop" i chcemy oszacować wartość stanów przy tej strategii. Metody: Monte Carlo prediction, TD(0), dynamic programming, jeśli model jest znany.

\subsubsection{Aktywne uczenie ze wzmocnieniem}

W aktywnym RL agent ma znaleźć dobrą albo optymalną politykę. Musi jednocześnie:
\begin{itemize}
  \item eksplorować, aby poznać skutki akcji,
  \item eksploatować, aby wybierać akcje już znane jako dobre.
\end{itemize}

To problem exploration-exploitation. Strategie:
\begin{itemize}
  \item $\eps$-greedy - z prawdopodobieństwem $1-\eps$ wybierz najlepszą znaną akcję, z $\eps$ losową,
  \item softmax/Boltzmann exploration,
  \item upper confidence bound,
  \item Thompson sampling,
  \item dodawanie szumu do parametrów albo akcji.
\end{itemize}

\subsubsection{On-policy i off-policy}

W uczeniu on-policy uczymy wartość albo poprawiamy tę samą politykę, która generuje dane. Przykład: SARSA.

W uczeniu off-policy dane generuje polityka behawioralna $b$, a uczymy politykę docelową $\pi$. Przykład: Q-learning. Off-policy pozwala uczyć politykę zachłanną na danych z polityki eksploracyjnej, ale może być mniej stabilne, zwłaszcza z aproksymacją funkcji.

\subsubsection{Funkcje wartości}

Wartość stanu:
\[
V^\pi(s)=\mathbb{E}_\pi[G_t\mid S_t=s].
\]
Wartość akcji:
\[
Q^\pi(s,a)=\mathbb{E}_\pi[G_t\mid S_t=s,A_t=a].
\]
Funkcja $V$ mówi, jak dobry jest stan przy danej polityce. Funkcja $Q$ mówi, jak dobra jest akcja w stanie, dlatego jest szczególnie użyteczna, gdy nie znamy modelu przejść.

\subsubsection{Co powiedzieć na obronie}

RL polega na uczeniu przez interakcję: stan, akcja, nagroda, kolejny stan. Polityka mapuje stany na akcje albo rozkłady akcji. Pasywne RL ocenia zadaną politykę, aktywne RL szuka polityki optymalnej i musi rozwiązać dylemat eksploracji i eksploatacji. Wartości $V^\pi$ i $Q^\pi$ opisują oczekiwany zwrot.

\subsection{Decyzyjny proces Markowa: definicja i przykłady}

MDP jest formalnym modelem sekwencyjnego podejmowania decyzji przy niepewności. Zakłada własność Markowa: przyszłość zależy od bieżącego stanu i akcji, a nie od pełnej historii.

\subsubsection{Definicja}

Skończony MDP można zapisać jako
\[
M=(S,A,p,r,\gamma),
\]
gdzie:
\begin{itemize}
  \item $S$ - zbiór stanów,
  \item $A$ - zbiór akcji,
  \item $p(s',r|s,a)$ - prawdopodobieństwo przejścia do stanu $s'$ i otrzymania nagrody $r$ po wykonaniu akcji $a$ w stanie $s$,
  \item $r$ - funkcja nagrody, często jako wartość oczekiwana,
  \item $\gamma\in[0,1]$ - współczynnik dyskontowy.
\end{itemize}

Warunek normalizacji:
\[
\sum_{s'\in S}\sum_{r\in R}p(s',r|s,a)=1
\quad\text{dla każdego }s,a.
\]

Z rozkładu $p(s',r|s,a)$ można wyznaczyć prawdopodobieństwo przejścia:
\[
p(s'|s,a)=\sum_r p(s',r|s,a),
\]
oraz oczekiwaną nagrodę:
\[
r(s,a)=\sum_r\sum_{s'} r\,p(s',r|s,a).
\]

\subsubsection{Własność Markowa}

Proces spełnia własność Markowa, jeśli:
\[
P(S_{t+1},R_{t+1}\mid S_0,A_0,\dots,S_t,A_t)
=P(S_{t+1},R_{t+1}\mid S_t,A_t).
\]
Stan powinien więc zawierać całą informację z historii potrzebną do przewidywania przyszłości. Jeśli obserwacja nie jest pełnym stanem, mamy POMDP, czyli częściowo obserwowalny MDP.

\subsubsection{Równania Bellmana}

Dla ustalonej polityki:
\[
V^\pi(s)=\sum_a\pi(a|s)\sum_{s',r}p(s',r|s,a)\left[r+\gamma V^\pi(s')\right].
\]
Dla funkcji akcji:
\[
Q^\pi(s,a)=\sum_{s',r}p(s',r|s,a)\left[r+\gamma\sum_{a'}\pi(a'|s')Q^\pi(s',a')\right].
\]

Optymalna funkcja wartości spełnia:
\[
V^*(s)=\max_a\sum_{s',r}p(s',r|s,a)\left[r+\gamma V^*(s')\right].
\]
Optymalna funkcja akcji:
\[
Q^*(s,a)=\sum_{s',r}p(s',r|s,a)\left[r+\gamma\max_{a'}Q^*(s',a')\right].
\]

\subsubsection{Przykłady MDP}

\begin{description}
  \item[Recycling robot] stan to poziom baterii, np. high/low. Akcje to search, wait, recharge. Nagrody zależą od zebrania śmieci i kar za rozładowanie.
  \item[Gridworld] stan to pozycja agenta na kratownicy, akcje to ruchy, przejścia mogą być stochastyczne, a nagrody przypisane do pól terminalnych i kosztów ruchu.
  \item[Blackjack] stan opisuje sumę kart gracza, odkrytą kartę krupiera i posiadanie używalnego asa. Akcje to hit/stand, nagroda końcowa to wygrana, przegrana albo remis.
  \item[Cart-pole] stan obejmuje pozycję wózka, prędkość, kąt słupka i prędkość kątową. Akcje to siła w lewo/prawo, nagroda za utrzymanie równowagi.
  \item[Bioreaktor] stan to odczyty sensorów i skład, akcje to ustawienia temperatury i mieszania, nagroda to tempo produkcji pożądanej substancji.
\end{description}

\subsubsection{Planowanie i uczenie}

Jeśli model $p$ i $r$ jest znany, możemy planować metodami programowania dynamicznego: policy evaluation, policy iteration, value iteration. Jeśli model nie jest znany, agent uczy się z doświadczenia metodami Monte Carlo, TD, SARSA, Q-learning albo metodami z aproksymacją funkcji.

\subsubsection{Co powiedzieć na obronie}

MDP to model $(S,A,p,r,\gamma)$, w którym przyszłość zależy tylko od bieżącego stanu i akcji. Polityka wybiera akcje, funkcje $V^\pi$ i $Q^\pi$ opisują oczekiwany zwrot, a równania Bellmana wiążą wartość stanu z nagrodą i wartością następnych stanów. Przykłady to robot sprzątający, gridworld, blackjack, cart-pole i sterowanie procesem.

\subsection{Metoda różnic czasowych TD-learning: idea, zakres stosowania, wady i zalety oraz porównanie z Monte Carlo}

Temporal Difference learning łączy idee Monte Carlo i programowania dynamicznego. Tak jak Monte Carlo uczy się z doświadczenia, bez pełnego modelu środowiska. Tak jak programowanie dynamiczne używa bootstrappingu, czyli aktualizuje wartość na podstawie obecnego oszacowania wartości następnego stanu.

\subsubsection{TD(0)}

Dla ustalonej polityki $\pi$ aktualizacja TD(0) ma postać:
\[
V(S_t)\leftarrow V(S_t)+\alpha\left[R_{t+1}+\gamma V(S_{t+1})-V(S_t)\right].
\]
Wyrażenie
\[
\delta_t=R_{t+1}+\gamma V(S_{t+1})-V(S_t)
\]
to błąd TD.

Cel aktualizacji TD to:
\[
R_{t+1}+\gamma V(S_{t+1}),
\]
czyli jednkrokowe przybliżenie zwrotu. W Monte Carlo celem jest pełny zwrot $G_t$ po zakończeniu epizodu.

\subsubsection{Algorytm predykcji TD(0)}

\begin{enumerate}
  \item Zainicjalizuj $V(s)$ dowolnie.
  \item Dla każdego epizodu:
  \begin{enumerate}
    \item rozpocznij w stanie $S$,
    \item wybierz akcję zgodnie z polityką $\pi$,
    \item wykonaj akcję, obserwuj $R$ i $S'$,
    \item zaktualizuj
    \[
    V(S)\leftarrow V(S)+\alpha\left[R+\gamma V(S')-V(S)\right],
    \]
    \item ustaw $S\leftarrow S'$,
    \item powtarzaj do stanu terminalnego.
  \end{enumerate}
\end{enumerate}

\subsubsection{Porównanie TD i Monte Carlo}

\begin{longtable}{p{0.24\textwidth}p{0.32\textwidth}p{0.32\textwidth}}
\toprule
Cecha & Monte Carlo & TD \\
\midrule
Model środowiska & nie wymaga & nie wymaga \\
Moment aktualizacji & po epizodzie & po każdym kroku \\
Cel aktualizacji & pełny zwrot $G_t$ & $R_{t+1}+\gamma V(S_{t+1})$ \\
Epizodyczność & wymaga epizodów albo specjalnych wariantów & działa w zadaniach ciągłych \\
Bootstrapping & nie & tak \\
Bias & mniejszy przy poprawnych próbkach & większy przez bootstrapping \\
Wariancja & często większa & zwykle mniejsza \\
Szybkość uczenia online & wolniejsza informacja & szybka aktualizacja \\
\bottomrule
\end{longtable}

MC jest nieobciążone względem rzeczywistych zwrotów przy próbkowaniu z polityki, ale może mieć dużą wariancję. TD wprowadza bias, bo używa własnych oszacowań, ale zwykle ma mniejszą wariancję i uczy się szybciej w zadaniach sekwencyjnych.

\subsubsection{SARSA}

SARSA jest on-policy metodą kontroli TD:
\[
Q(S_t,A_t)\leftarrow Q(S_t,A_t)+\alpha\left[R_{t+1}+\gamma Q(S_{t+1},A_{t+1})-Q(S_t,A_t)\right].
\]
Nazwa pochodzi od sekwencji:
\[
S_t,A_t,R_{t+1},S_{t+1},A_{t+1}.
\]
Ponieważ używa akcji faktycznie wybranej przez bieżącą politykę, uczy wartości polityki eksploracyjnej.

\subsubsection{Q-learning}

Q-learning jest off-policy:
\[
Q(S_t,A_t)\leftarrow Q(S_t,A_t)+\alpha\left[R_{t+1}+\gamma\max_a Q(S_{t+1},a)-Q(S_t,A_t)\right].
\]
Uczy wartości polityki zachłannej, nawet jeśli dane są generowane przez politykę eksploracyjną, np. $\eps$-greedy.

\subsubsection{Zakres stosowania}

TD jest użyteczne, gdy:
\begin{itemize}
  \item środowisko jest nieznane,
  \item dane przychodzą online,
  \item epizody są długie albo nie kończą się,
  \item chcemy aktualizować po każdym kroku,
  \item pełny zwrot Monte Carlo miałby dużą wariancję.
\end{itemize}

TD jest podstawą wielu metod RL, także z aproksymacją funkcji i deep RL.

\subsubsection{Wady}

\begin{itemize}
  \item bias przez bootstrapping,
  \item wrażliwość na krok uczenia $\alpha$,
  \item możliwa niestabilność przy aproksymacji funkcji, off-policy i bootstrappingu,
  \item potrzeba eksploracji dla metod kontroli,
  \item trudność doboru reprezentacji stanu.
\end{itemize}

\subsubsection{Zalety}

\begin{itemize}
  \item uczenie online,
  \item brak potrzeby modelu środowiska,
  \item brak potrzeby czekania do końca epizodu,
  \item często mniejsza wariancja niż MC,
  \item możliwość stosowania w zadaniach ciągłych,
  \item naturalne połączenie z funkcjami wartości akcji.
\end{itemize}

\subsubsection{Co powiedzieć na obronie}

TD-learning aktualizuje wartość stanu po każdym kroku, używając różnicy między starą wartością a celem $R+\gamma V(S')$. Błąd TD to $\delta=R+\gamma V(S')-V(S)$. W porównaniu z Monte Carlo TD nie czeka do końca epizodu i zwykle ma mniejszą wariancję, ale jest obciążone przez bootstrapping. SARSA jest on-policy, Q-learning off-policy.

\section{Sieci neuronowe}

\subsection{Uczenie nadzorowane i nienadzorowane: definicje i przykłady}

Uczenie sieci neuronowej polega na zmianie wag połączeń tak, aby sieć realizowała pożądaną funkcję. Struktura połączeń i wartości wag określają, jakie przekształcenie wejścia w wyjście wykonuje sieć. W klasycznym podziale wyróżnia się uczenie nadzorowane, nienadzorowane i uczenie ze wzmocnieniem.

\subsubsection{Uczenie nadzorowane}

W uczeniu nadzorowanym dane treningowe mają postać par:
\[
(x^{(i)},y^{(i)}),\qquad i=1,\dots,n,
\]
gdzie $x^{(i)}\in\R^m$ jest wektorem wejściowym, a $y^{(i)}\in\R^k$ pożądanym wyjściem. Zadaniem jest nauczenie funkcji:
\[
f_\theta:\R^m\to\R^k
\]
minimalizującej błąd predykcji na nowych danych.

Uczenie polega na minimalizacji funkcji straty:
\[
\min_\theta \frac{1}{n}\sum_{i=1}^{n} \ell(f_\theta(x^{(i)}),y^{(i)}).
\]
Dla regresji często stosuje się błąd kwadratowy:
\[
L=\frac{1}{2}\sum_i\sum_j(y_j^{(i)}-\hat y_j^{(i)})^2.
\]
Dla klasyfikacji wieloklasowej z softmaxem typowa jest entropia krzyżowa:
\[
CE=-\sum_{j=1}^{k}y_j\log \hat y_j.
\]
Przy one-hot encoding sprowadza się to do $-\log$ prawdopodobieństwa klasy prawdziwej.

Przykłady:
\begin{itemize}
  \item klasyfikacja obrazów,
  \item rozpoznawanie cyfr,
  \item klasyfikacja spamu,
  \item regresja cen mieszkań,
  \item prognoza popytu,
  \item segmentacja obrazu z etykietami pikseli.
\end{itemize}

\subsubsection{Uczenie nienadzorowane}

W uczeniu nienadzorowanym mamy tylko wejścia:
\[
x^{(1)},\dots,x^{(n)}
\]
bez etykiet. Celem jest odkrycie struktury danych: grup, reprezentacji, czynników zmienności albo rozkładu.

Przykłady zadań:
\begin{itemize}
  \item klastrowanie,
  \item redukcja wymiaru,
  \item detekcja anomalii,
  \item kompresja reprezentacji,
  \item modelowanie generatywne,
  \item samoorganizujące mapy cech.
\end{itemize}

W sieciach neuronowych przykładami są autoenkodery, Restricted Boltzmann Machines, sieci Kohonena, uczenie Hebbowskie, WTA/WTM oraz modele generatywne.

\subsubsection{WTA i WTM}

Winner Takes All oznacza, że dla danego wejścia zwycięża neuron o największej aktywacji, a uczeniu podlega głównie on. Winner Takes Most rozszerza to na sąsiedztwo zwycięzcy. WTM jest podstawą samoorganizujących map Kohonena: aktualizuje się zwycięski neuron i jego sąsiadów, co prowadzi do topologicznego uporządkowania mapy.

\subsubsection{Porównanie}

\begin{longtable}{p{0.24\textwidth}p{0.32\textwidth}p{0.32\textwidth}}
\toprule
Cecha & Uczenie nadzorowane & Uczenie nienadzorowane \\
\midrule
Dane & wejścia i etykiety & same wejścia \\
Cel & predykcja znanego celu & odkrycie struktury \\
Funkcja straty & porównuje predykcję z etykietą & rekonstrukcja, podobieństwo, gęstość, topologia \\
Przykłady & klasyfikacja, regresja & klastrowanie, PCA, SOM, autoenkodery \\
Ocena & accuracy, F1, MSE & silhouette, rekonstrukcja, ocena jakości reprezentacji \\
\bottomrule
\end{longtable}

\subsubsection{Co powiedzieć na obronie}

Uczenie nadzorowane korzysta z par wejście-wyjście i minimalizuje błąd względem etykiet. Uczenie nienadzorowane nie ma etykiet i szuka struktury w danych, np. klastrów albo reprezentacji. Przykładem nadzorowanym jest MLP klasyfikujący cyfry, a nienadzorowanym mapa Kohonena albo autoenkoder.

\subsection{Porównanie modeli sieciowych Hopfielda, Grossberga i Kohonena}

Modele Hopfielda, Grossberga i Kohonena należą do klasycznych architektur sieci neuronowych, ale rozwiązują inne zadania i mają inne mechanizmy uczenia.

\subsubsection{Sieć Hopfielda}

Sieć Hopfielda jest rekurencyjną siecią z pełnymi symetrycznymi połączeniami między neuronami i bez połączeń własnych:
\[
w_{ij}=w_{ji},\qquad w_{ii}=0.
\]
Klasyczny model dyskretny ma neurony bipolarne $v_i\in\{-1,1\}$ albo unipolarne $v_i\in\{0,1\}$. Stan sieci jest wektorem aktywacji. Aktualizacja może być synchroniczna albo asynchroniczna.

Główne zastosowanie: pamięć skojarzeniowa, czyli content-addressable memory. Sieć przechowuje wzorce jako atraktory dynamiki. Po podaniu wzorca zaszumionego powinna zbiec do najbliższego zapamiętanego wzorca.

\subsubsection{Modele Grossberga}

Stephen Grossberg rozwijał modele konkurencyjnego i adaptacyjnego uczenia, w tym Adaptive Resonance Theory. W kontekście klasycznego porównania z Hopfieldem i Kohonenem Grossberg kojarzy się z sieciami:
\begin{itemize}
  \item uczącymi się stabilnie nowych kategorii,
  \item opartymi na konkurencji neuronów,
  \item rozwiązującymi dylemat stability-plasticity,
  \item używającymi mechanizmu rezonansu między reprezentacją wejścia i kategorią.
\end{itemize}

ART ma warstwę wejściową i warstwę kategorii. Wejście aktywuje kategorię kandydującą; jeśli podobieństwo przekracza parametr czujności, następuje rezonans i uczenie. Jeśli nie, kategoria jest odrzucana i szukana jest inna albo tworzona nowa. Wysoka czujność tworzy więcej, bardziej szczegółowych kategorii; niska czujność daje kategorie ogólniejsze.

Główne zastosowania: klastrowanie, klasyfikacja inkrementalna, uczenie bez katastroficznego nadpisywania, adaptacyjne rozpoznawanie wzorców.

\subsubsection{Sieć Kohonena}

Self-Organizing Map jest nienadzorowaną siecią konkurencyjną. Neurony są ułożone zwykle na siatce 1D albo 2D, a każdy neuron ma wektor wag $w_i$ w przestrzeni danych. Dla wejścia $x$ wybiera się best matching unit:
\[
c=\argmin_i\|x-w_i\|.
\]
Następnie aktualizuje się zwycięzcę i jego sąsiadów:
\[
w_i(t+1)=w_i(t)+\alpha(t)h_{ci}(t)(x-w_i(t)),
\]
gdzie $h_{ci}$ maleje z odległością neuronu $i$ od zwycięzcy $c$ na mapie.

SOM tworzy odwzorowanie danych wysokowymiarowych na niskowymiarową siatkę z zachowaniem lokalnej topologii. Służy do wizualizacji, klastrowania i eksploracji danych.

\subsubsection{Porównanie}

\begin{longtable}{p{0.2\textwidth}p{0.24\textwidth}p{0.24\textwidth}p{0.24\textwidth}}
\toprule
Cecha & Hopfield & Grossberg/ART & Kohonen/SOM \\
\midrule
Typ & rekurencyjna pamięć skojarzeniowa & adaptacyjne kategorie, konkurencja & samoorganizująca mapa \\
Uczenie & często Hebbowskie albo jego modyfikacje & rezonans, dopasowanie, czujność & WTA/WTM, uczenie sąsiedztwa \\
Zadanie & odtwarzanie wzorców, optymalizacja & stabilne uczenie kategorii & klastrowanie i wizualizacja \\
Dynamika & zbieganie do atraktora & wybór lub tworzenie kategorii & porządkowanie topologiczne \\
Nadzór & zwykle nienadzorowane/asocjacyjne & nienadzorowane lub nadzorowane warianty ART & nienadzorowane \\
Problem kluczowy & pojemność i minima fałszywe & stability-plasticity & dobór mapy, sąsiedztwa i metryki \\
\bottomrule
\end{longtable}

\subsubsection{Co powiedzieć na obronie}

Hopfield to rekurencyjna sieć pamięci skojarzeniowej, w której wzorce są atraktorami. Grossberg/ART koncentruje się na adaptacyjnym tworzeniu kategorii i kompromisie stabilność-plastyczność. Kohonen/SOM to nienadzorowana mapa samoorganizująca, która zachowuje topologię danych i nadaje się do klastrowania oraz wizualizacji.

\subsection{Dobór architektury sieci neuronowej do wybranego zadania}

Dobór architektury powinien wynikać z rodzaju danych, zadania, ograniczeń obliczeniowych i wymagań jakościowych. Nie istnieje jedna najlepsza sieć do wszystkiego.

\subsubsection{Klasyfikacja tabularna}

Dla danych tabelarycznych często warto zacząć od prostych modeli: regresji logistycznej, drzew, random forest, gradient boosting. Jeśli używamy sieci, typowy jest MLP:
\[
x\to \text{Dense}\to \text{aktywacja}\to \dots \to \text{softmax/sigmoid}.
\]
Dla klasyfikacji wieloklasowej ostatnia warstwa ma $k$ neuronów i softmax, a strata to categorical cross-entropy. Dla binarnej - jeden neuron z sigmoidą i binary cross-entropy.

\subsubsection{Regresja}

Dla regresji używa się MLP z liniowym neuronem wyjściowym albo odpowiednią aktywacją zgodną z zakresem wartości. Strata: MSE, MAE, Huber. Jeśli predykcja ma być nieujemna, można użyć softplus albo modelować logarytm zmiennej.

\subsubsection{Obrazy}

Dla obrazów naturalnym wyborem są sieci konwolucyjne, bo wykorzystują lokalność i współdzielenie wag. Warstwy konwolucyjne wykrywają lokalne wzorce, głębsze warstwy budują cechy złożone. Dla segmentacji stosuje się architektury encoder-decoder, np. U-Net. Dla klasyfikacji można użyć transfer learningu z modelu wstępnie trenowanego.

\subsubsection{Sekwencje i szeregi czasowe}

Możliwości:
\begin{itemize}
  \item RNN, LSTM, GRU - klasyczne modele rekurencyjne,
  \item 1D CNN - lokalne wzorce w czasie,
  \item Transformery - długie zależności i równoległe przetwarzanie,
  \item modele temporal convolutional networks,
  \item architektury hybrydowe.
\end{itemize}

Dla predykcji szeregów czasowych ważne są okna czasowe, cechy kalendarzowe, unikanie przecieku przyszłości i walidacja krocząca.

\subsubsection{Tekst}

Współcześnie standardem są transformery i embeddingi. Dla prostych zadań można użyć bag-of-words z modelem liniowym albo małego MLP. Dla zadań semantycznych używa się modeli językowych, fine-tuningu albo klasyfikacji na embeddingach.

\subsubsection{Uczenie nienadzorowane i kompresja}

Dla kompresji i reprezentacji stosuje się autoenkodery. Bottleneck wymusza kod niskowymiarowy. Dla wizualizacji klasyczne są SOM, PCA, t-SNE, UMAP. Dla danych obrazowych autoenkoder konwolucyjny zachowuje strukturę przestrzenną.

\subsubsection{Pamięć skojarzeniowa i odtwarzanie wzorców}

Dla zadania "po zaszumionym wzorcu odtwórz najbliższy zapamiętany wzorzec" naturalna jest sieć Hopfielda. Dla dużych i złożonych wzorców klasyczny Hopfield ma ograniczoną pojemność, więc stosuje się modyfikacje reguł uczenia albo nowocześniejsze pamięci asocjacyjne.

\subsubsection{Kryteria doboru rozmiaru}

Za mała sieć ma duży bias i underfitting. Za duża może mieć overfitting, choć przy dużych danych i regularyzacji duże sieci mogą generalizować dobrze. Dobór obejmuje:
\begin{itemize}
  \item liczbę warstw,
  \item liczbę neuronów/kanałów,
  \item funkcje aktywacji,
  \item normalizację,
  \item dropout,
  \item regularyzację wag,
  \item learning rate i optimizer,
  \item batch size,
  \item early stopping.
\end{itemize}

\subsubsection{Walidacja architektury}

Architekturę dobiera się eksperymentalnie na zbiorze walidacyjnym. Należy porównywać z baseline'ami, analizować krzywe uczenia, sprawdzać overfitting i wykonywać ablation study. Dla małych danych często lepszy jest transfer learning albo prostszy model.

\subsubsection{Co powiedzieć na obronie}

Architekturę dobiera się do struktury danych i zadania: MLP dla danych wektorowych, CNN dla obrazów, RNN/LSTM/Transformer dla sekwencji, autoenkoder dla kompresji, SOM dla mapowania nienadzorowanego, Hopfield dla pamięci skojarzeniowej. Dobór rozmiaru i regularyzacji trzeba potwierdzić walidacją.

\subsection{Idea i własności pamięci skojarzeniowych Hopfielda}

Pamięć skojarzeniowa Hopfielda adresuje zawartością, nie adresem. Oznacza to, że po podaniu niepełnego albo zaszumionego wzorca sieć powinna odtworzyć najbliższy zapamiętany wzorzec.

\subsubsection{Architektura}

Klasyczna dyskretna sieć Hopfielda:
\begin{itemize}
  \item ma $N$ neuronów,
  \item każdy neuron jest połączony z każdym innym,
  \item wagi są symetryczne: $w_{ij}=w_{ji}$,
  \item brak połączeń własnych: $w_{ii}=0$,
  \item stan neuronu jest bipolarny $v_i\in\{-1,1\}$.
\end{itemize}

Wejście neuronu:
\[
u_i=\sum_{j=1}^{N}w_{ij}v_j-\theta_i,
\]
a aktualizacja:
\[
v_i\leftarrow \operatorname{sgn}(u_i).
\]

\subsubsection{Uczenie Hebbowskie}

Dla $M$ wzorców bipolarnych $x^1,\dots,x^M$, $x^s\in\{-1,1\}^N$, klasyczna reguła Hebba:
\[
w_{ij}=
\begin{cases}
0, & i=j,\\
\frac{1}{N}\sum_{s=1}^{M}x_i^s x_j^s, & i\ne j.
\end{cases}
\]
Macierzowo:
\[
W=\frac{1}{N}(XX^T-MI)
\]
przy odpowiedniej konwencji zapisu wzorców.

Reguła wzmacnia połączenia neuronów aktywnych razem i osłabia połączenia aktywnych przeciwnie.

\subsubsection{Faza rozpoznawania}

Podajemy wektor startowy $z$, np. zaszumiony wzorzec. Sieć iteracyjnie aktualizuje neurony. Możliwe wyniki:
\begin{itemize}
  \item zbieżność do zapamiętanego wzorca,
  \item zbieżność do wzorca fałszywego,
  \item w trybie synchronicznym cykl długości 2,
  \item błędne odtworzenie przy zbyt dużym szumie albo przeciążeniu pamięci.
\end{itemize}

\subsubsection{Funkcja energii}

Dla symetrycznych wag i braku wag własnych definiuje się energię:
\[
E(v)=-\frac{1}{2}\sum_{i\ne j}w_{ij}v_iv_j+\sum_i\theta_i v_i.
\]
Przy aktualizacji asynchronicznej energia nie rośnie. Ponieważ liczba stanów jest skończona, sieć zbiega do punktu stałego. To podstawowa własność stabilizująca model.

\subsubsection{Atraktory i baseny przyciągania}

Zapamiętane wzorce powinny być minimami lokalnymi energii. Basen przyciągania wzorca to zbiór stanów startowych, z których dynamika zbiega do tego wzorca. Im większy basen, tym większa odporność na szum. Wzorce zbyt podobne do siebie albo zbyt liczne zmniejszają baseny i zwiększają liczbę atraktorów fałszywych.

\subsubsection{Pojemność}

Dla losowych niezależnych wzorców bipolarnych klasyczna praktyczna estymacja pojemności reguły Hebba wynosi około:
\[
M\approx 0.15N.
\]
Asymptotyczne wyniki często podają skalowanie rzędu:
\[
M\sim \frac{N}{4\ln N}
\]
dla określonych warunków bezbłędnego odtwarzania. Pojemność zależy od definicji dopuszczalnego błędu, korelacji wzorców i reguły uczenia.

\subsubsection{Modyfikacje}

Wady Hebba: ograniczona pojemność, wzorce fałszywe, problemy dla wzorców skorelowanych i rzadkich. Modyfikacje:
\begin{itemize}
  \item reguła pseudoodwrotna/projection rule,
  \item iteracyjne uczenie macierzy wag,
  \item dobór progów indywidualnych,
  \item preprocessing wzorców,
  \item uogólnione sieci Hopfielda.
\end{itemize}
Zwykle poprawa pojemności kosztuje utratę lokalności albo prostoty biologicznej interpretacji.

\subsubsection{Co powiedzieć na obronie}

Sieć Hopfielda jest rekurencyjną pamięcią skojarzeniową. Wzorce są zapisywane w wagach, często regułą Hebba, i odpowiadają atraktorom funkcji energii. Po podaniu zaszumionego wzorca sieć przez iteracyjne aktualizacje zbiega do lokalnego minimum energii. W trybie asynchronicznym przy wagach symetrycznych i zerowej przekątnej energia nie rośnie, więc sieć zbiega. Ograniczenia to pojemność, atraktory fałszywe i wrażliwość na korelacje wzorców.

\subsection{Algorytm propagacji wstecznej błędu, własności i podstawowe modyfikacje}

Backpropagation jest podstawowym algorytmem obliczania gradientów w wielowarstwowych sieciach neuronowych. Łączy przejście w przód, obliczenie straty i propagację gradientu od wyjścia do wejścia za pomocą reguły łańcuchowej.

\subsubsection{Przejście w przód}

Dla warstwy $l$:
\[
z^{[l]}=W^{[l]}a^{[l-1]}+b^{[l]},
\]
\[
a^{[l]}=g^{[l]}(z^{[l]}).
\]
Na końcu otrzymujemy predykcję $\hat y=a^{[L]}$ i stratę $L(\hat y,y)$.

\subsubsection{Przejście wstecz}

Backpropagation oblicza pochodne straty względem parametrów. Dla warstwy wyjściowej wyznacza się błąd $\delta^{[L]}$, a następnie dla warstw wcześniejszych:
\[
\delta^{[l]}=\left((W^{[l+1]})^T\delta^{[l+1]}\right)\odot g'^{[l]}(z^{[l]}).
\]
Gradienty:
\[
\frac{\partial L}{\partial W^{[l]}}=\delta^{[l]}(a^{[l-1]})^T,\qquad
\frac{\partial L}{\partial b^{[l]}}=\delta^{[l]}.
\]
Wagi aktualizuje się np. gradient descent:
\[
W\leftarrow W-\eta\nabla_W L.
\]

\subsubsection{Tryby uczenia}

\begin{description}
  \item[Online/stochastic] aktualizacja po pojedynczym przykładzie.
  \item[Mini-batch] aktualizacja po małej paczce przykładów; standard praktyczny.
  \item[Batch] aktualizacja po całym zbiorze treningowym.
\end{description}

Epoka oznacza jedno przejście przez cały zbiór treningowy. Iteracja to pojedyncza aktualizacja parametrów.

\subsubsection{Warunki stopu}

\begin{itemize}
  \item maksymalna liczba epok,
  \item strata poniżej progu,
  \item norma aktualizacji poniżej progu,
  \item brak poprawy na walidacji,
  \item early stopping.
\end{itemize}

\subsubsection{Zalety backpropagation}

\begin{itemize}
  \item uniwersalność dla różniczkowalnych architektur,
  \item efektywne obliczanie gradientów przez reuse wyników pośrednich,
  \item możliwość trenowania MLP i głębokich sieci,
  \item zgodność z wieloma funkcjami strat,
  \item łatwe połączenie z optymalizatorami.
\end{itemize}

\subsubsection{Problemy}

Z wykładów istotne problemy:
\begin{itemize}
  \item minima lokalne i punkty siodłowe,
  \item wolna zbieżność w płaskich obszarach funkcji błędu,
  \item oscylacje przy zbyt dużym kroku,
  \item koszt obliczeniowy,
  \item zanikające i eksplodujące gradienty,
  \item catastrophic interference/forgetting,
  \item dobór architektury.
\end{itemize}

\subsubsection{Momentum}

Momentum dodaje bezwładność:
\[
\Delta w(t)=-\eta\nabla E(w(t))+\alpha\Delta w(t-1).
\]
Pomaga ograniczać oscylacje i przyspieszać ruch w długich płaskich dolinach funkcji błędu. Zbyt duże momentum może powodować niestabilność.

\subsubsection{Adaptacyjny learning rate}

Jeśli znak pochodnej cząstkowej dla danej wagi pozostaje taki sam, można zwiększać krok; jeśli znak się zmienia, prawdopodobnie przeskoczyliśmy minimum i krok trzeba zmniejszyć. To idea metod takich jak Silva-Almeida, Delta-bar-delta i Rprop.

Rprop używa głównie znaku gradientu, a nie jego wartości:
\[
\Delta w_i=-\eta_i \operatorname{sgn}\left(\frac{\partial E}{\partial w_i}\right),
\]
przy adaptacji $\eta_i$ zależnej od zmian znaku gradientu. Jest metodą batchową i dobrze radzi sobie z płaskimi regionami.

\subsubsection{QuickProp i metody drugiego rzędu}

Metody drugiego rzędu wykorzystują krzywiznę funkcji błędu. Newton:
\[
w_{k+1}=w_k-\left(\nabla^2E(w_k)\right)^{-1}\nabla E(w_k).
\]
Jest szybki lokalnie, ale Hessian i jego odwrotność są kosztowne. QuickProp aproksymuje informację drugiego rzędu dla pojedynczych wag i może przyspieszać uczenie, ale wymaga ostrożności.

\subsubsection{Regularyzacja i stabilizacja}

Podstawowe modyfikacje współczesne:
\begin{itemize}
  \item weight decay,
  \item dropout,
  \item batch normalization/layer normalization,
  \item data augmentation,
  \item early stopping,
  \item gradient clipping,
  \item lepsze inicjalizacje, np. Xavier/He,
  \item ReLU i jej warianty zamiast sigmoid w głębokich sieciach,
  \item Adam/AdamW.
\end{itemize}

\subsubsection{Co powiedzieć na obronie}

Backpropagation oblicza gradient funkcji straty względem wag przez przejście w przód i propagację błędu wstecz regułą łańcuchową. Wagi aktualizuje się gradientowo. Zalety to efektywność i uniwersalność; problemy to minima lokalne, wolna zbieżność, oscylacje, zanik gradientu i katastroficzne zapominanie. Modyfikacje obejmują momentum, adaptacyjne kroki, Rprop, QuickProp, regularyzację, dropout i early stopping.

\subsection{Problem katastroficznego zapominania: na czym polega i jak przeciwdziałać}

Katastroficzne zapominanie, nazywane też catastrophic interference, polega na tym, że sieć uczona sekwencyjnie nowych zadań albo klas traci wcześniej nabytą wiedzę. Nowe aktualizacje wag nadpisują reprezentacje potrzebne dla starych zadań.

\subsubsection{Dlaczego występuje}

W klasycznym uczeniu backpropagation wszystkie zadania współdzielą te same parametry. Jeśli uczymy najpierw zadania A, a potem tylko zadania B, gradient z B zmienia wagi bez informacji o tym, że część wag była istotna dla A. W efekcie spada jakość na A.

Problem jest szczególnie silny, gdy:
\begin{itemize}
  \item dane przychodzą sekwencyjnie,
  \item nie można przechowywać starych danych,
  \item zadania korzystają z tych samych parametrów,
  \item rozkłady danych zmieniają się,
  \item model ma ograniczoną pojemność albo zbyt agresywne uczenie.
\end{itemize}

\subsubsection{Stability-plasticity dilemma}

Model powinien być plastyczny, aby uczyć się nowych informacji, ale stabilny, aby nie tracić starych. Zbyt duża plastyczność powoduje zapominanie. Zbyt duża stabilność blokuje uczenie nowych zadań.

\subsubsection{Interleaved learning i rehearsal}

Najprostsza metoda to mieszać stare i nowe przykłady. Rehearsal przechowuje bufor reprezentatywnych starych danych i podczas uczenia nowych zadań powtarza stare. Warianty:
\begin{itemize}
  \item experience replay,
  \item exemplar memory,
  \item generative replay - starych danych dostarcza model generatywny,
  \item pseudo-rehearsal - generowanie syntetycznych przykładów.
\end{itemize}

Wada: potrzeba przechowywania danych albo generatora, koszt i kwestie prywatności.

\subsubsection{Regularyzacja ważnych wag}

Metody takie jak Elastic Weight Consolidation ograniczają zmianę wag ważnych dla poprzednich zadań:
\[
L(\theta)=L_{new}(\theta)+\lambda\sum_i F_i(\theta_i-\theta_i^*)^2,
\]
gdzie $F_i$ mierzy ważność parametru, np. przez informację Fishera. Podobną ideę mają Synaptic Intelligence i Memory Aware Synapses.

\subsubsection{Metody modularne}

Można przydzielać różne moduły albo części sieci do różnych zadań:
\begin{itemize}
  \item osobne głowy wyjściowe,
  \item progressive networks,
  \item adaptery,
  \item maski parametrów,
  \item dynamiczne rozszerzanie architektury.
\end{itemize}

Zaletą jest ochrona starej wiedzy. Wadą jest rosnący model i problem łączenia wyników modułów.

\subsubsection{Zamrażanie i współdzielenie cech}

W wykładach pojawia się idea incremental class learning: uczyć klasy kolejno, wyznaczać neurony reprezentujące istotne cechy klasy, zamrażać ich wagi i wykorzystywać je przy kolejnych klasach. Jeśli nowa klasa może użyć już zamrożonych cech, uczenie jest szybsze i mniej niszczy starą wiedzę. To podejście łączy stabilność zamrożonych reprezentacji z plastycznością niezamrożonych części sieci.

\subsubsection{Knowledge distillation}

Przy uczeniu nowego zadania można karać odchylenie odpowiedzi nowego modelu od odpowiedzi starego modelu na starych albo reprezentatywnych danych. Dzięki temu model zachowuje funkcję wejście-wyjście poprzedniej wersji.

\subsubsection{Architektury i reprezentacje}

Zapominanie można ograniczać przez:
\begin{itemize}
  \item większą pojemność modelu,
  \item reprezentacje disentangled,
  \item sparse activations,
  \item normalizację i małe kroki uczenia,
  \item task-specific adapters,
  \item metauczenie strategii ciągłego uczenia.
\end{itemize}

\subsubsection{Ocena ciągłego uczenia}

Trzeba mierzyć nie tylko wynik na ostatnim zadaniu, ale też:
\begin{itemize}
  \item średnią jakość na wszystkich zadaniach,
  \item forgetting measure,
  \item forward transfer - czy stare uczenie pomaga nowemu,
  \item backward transfer - czy nowe uczenie poprawia albo pogarsza stare,
  \item pamięć i koszt obliczeniowy.
\end{itemize}

\subsubsection{Co powiedzieć na obronie}

Katastroficzne zapominanie polega na utracie wcześniej nauczonej wiedzy podczas sekwencyjnego uczenia nowych zadań, bo te same wagi są modyfikowane przez nowe gradienty. Przeciwdziała się przez rehearsal, generative replay, regularyzację ważnych wag, modularność, zamrażanie i współdzielenie cech, dynamiczne rozszerzanie architektury oraz distillation. Sednem jest kompromis między stabilnością starej wiedzy i plastycznością potrzebną do uczenia nowej.


\end{document}
